
\medterm{L-Dopa} A medicine that the brain converts into dopamine. It is usually combined with carbidopa or benserazide to treat Parkinson disease.
\medterm{La Crosse Encephalitis} A mosquito-borne viral encephalitis caused by La Crosse virus, usually affecting children and causing fever, headache, seizures, or altered consciousness.
\medterm{LAAO} Alternate name; see \textit{Left Atrial Appendage Occlusion}.

\medterm{Labetalol} A medicine that lowers blood pressure by slowing the heart and relaxing blood vessels. It is also used for severe high blood pressure during pregnancy and may cause a slow heartbeat, low blood pressure, or breathing problems in people with asthma.

\medterm{Labia} The folds of skin surrounding the entrance to the vagina and urethra, consisting mainly of the labia majora and labia minora.

\medterm{Labia Majora} The paired outer folds of skin of the vulva that help protect the urethral and vaginal openings and contain connective tissue, glands, and fat.

\medterm{Labia Minora} The paired inner folds of tissue of the vulva, lying within the labia majora and surrounding the vestibule and its openings.

\medterm{Labial Sounds} Speech sounds made with one or both lips, such as the sounds in “p,” “b,” and “m.” Difficulty making them may occur with problems affecting the lips, mouth, nerves, or speech muscles.

\medterm{Labiaplasty} Surgical reshaping or reduction of the labia, usually the labia minora, performed for functional symptoms, reconstruction, or selected cosmetic indications. Risks include pain, bleeding, infection, scarring, and altered sensation.

\medterm{Labile} Easily or rapidly changing, such as blood pressure, blood sugar, mood, or symptoms. Labile measurements may reflect illness, medicines, stress, hormones, or changing treatment and often require repeated observation.

\medterm{Labile Hypertension} Blood pressure that changes substantially over short periods. Stress, pain, medicines, illness, and measurement conditions can contribute, so repeated properly performed readings help confirm the pattern.

\medterm{Lability} A tendency to change easily or unpredictably, such as rapidly shifting emotions or blood pressure. Its medical importance depends on the symptom, severity, and underlying cause.
\medterm{Labium} A lip-like anatomical structure; in obstetrics and gynecology, usually one of the labia of the vulva.
\medterm{Labium Majus} One of the paired outer folds of the vulva containing skin, connective tissue, glands, vessels, and nerves and protecting the more internal genital structures.

\medterm{Labium Minus} One of the paired inner folds of the vulva bordering the vestibule and helping protect the urethral and vaginal openings.

\medterm{Labor} The process in which regular contractions open the cervix and lead to delivery of the baby and placenta. It includes the stages of cervical opening, birth of the baby, and delivery of the placenta.

\medterm{Labor: Signs And Symptoms} Labor is the process of regular uterine contractions that opens the cervix and leads to birth. Possible signs include increasingly regular contractions, rupture of the membranes, a blood-stained mucus discharge, pelvic pressure, and back pain. Heavy bleeding, reduced fetal movement, severe pain, or suspected complications requires urgent assessment.

\medterm{Labor Complication} A maternal or fetal problem occurring during labor or delivery, such as abnormal labor progress, hemorrhage, infection, fetal distress, or malpresentation, requiring assessment and management based on the cause.

\medterm{Labor Dystocia} Labor that is unusually slow or difficult because of weak contractions, the baby's position or size, or the shape of the mother's pelvis. Treatment depends on the cause and may include medicines or an operation.

\medterm{Labor False (Braxton Hicks Contractions)} Alternate index label for Braxton Hicks Contractions.

\medterm{Labor Pain} Pain caused by uterine contractions, cervical dilation, and stretching of pelvic tissues during labor; its intensity and management vary with labor stage and individual factors.

\medterm{Laboratory} A facility or department where clinical specimens are examined using microbiologic, hematologic, chemical, molecular, or other diagnostic methods.

\medterm{Laboratory Medicine} The medical specialty that uses laboratory tests on blood, tissue, and other specimens for diagnosis, screening, monitoring, and treatment.
\medterm{Laboratory Procedure} A method used to collect, prepare, examine, or measure a patient sample such as blood, urine, tissue, or a swab. Results can help diagnose disease or guide treatment and must be interpreted with the clinical findings.

\medterm{Laboratory Test} An examination of a body fluid, tissue, or other specimen to obtain information for diagnosis, monitoring, screening, or treatment.

\medterm{Labour} The process of childbirth, consisting of regular uterine contractions that cause
progressive cervical effacement and dilation, leading to delivery of the baby and
placenta.

\medterm{Labra} The plural of labrum: tough rings of tissue around certain joint sockets, especially the shoulder and hip. They deepen the socket, improve stability, and help distribute pressure during movement.

\medterm{Labral Tear} A tear in the fibrocartilaginous rim of a joint socket, commonly the shoulder or hip, causing pain, clicking, instability, or reduced range of motion.

\medterm{Labrum} A ring of tough, flexible tissue around the socket of a joint. It deepens and stabilizes the shoulder or hip socket.

\medterm{Labyrinth} The inner-ear system that contains the cochlea for hearing and the balance organs that sense head movement and position. Disease or injury can cause dizziness, unsteadiness, ringing, or hearing loss.

\medterm{Labyrinth Disorder} A problem affecting inner-ear hearing or balance structures, potentially causing spinning sensations, nausea, unsteadiness, hearing changes, or abnormal eye movements.

\medterm{Labyrinthine Membrane} The delicate, fluid-containing system inside the inner ear that contains the organs for hearing and balance. Injury or abnormal fluid pressure can cause dizziness, hearing loss, or ringing in the ear.

\medterm{Labyrinthitis} Inflammation of the inner ear that can cause sudden spinning dizziness, imbalance, nausea, ringing in the ear, and sometimes hearing loss.

\medterm{Labyrinthitis (Inner Ear Inflammation)} Alternate index label for Inner Ear Inflammation.


\medterm{Laceration} A torn wound caused by blunt or irregular force. It may involve only the skin or extend into fat, muscle, tendons, nerves, or blood vessels; deep, gaping, dirty, or heavily bleeding wounds need medical care.

\medterm{Laceration (Forensic)} A torn wound caused by blunt force, usually with irregular bruised edges and tissue strands crossing the wound. These features help distinguish it from a clean cut caused by a sharp object.

\synonyms
Forensic

\medterm{Laceration - Sutures Or Staples - At Home} Home care after sutures or staples includes keeping the wound clean and protected, following dressing instructions, avoiding strain, and watching for infection or wound separation.

\medterm{Lacerations - Liquid Bandage} Liquid bandage is a film-forming adhesive applied over small, clean, shallow cuts to protect the wound; it should not be used on deep, infected, heavily bleeding, or gaping wounds.

\medterm{Lachman Test} A knee examination in which the lower leg is gently pulled forward while the knee is slightly bent. Too much movement may suggest an injury to the anterior cruciate ligament.

\medterm{Lachrymal} Relating to tears or the structures that produce and drain them; also spelled lacrimal.
\medterm{Lacosamide} An antiseizure medicine used mainly for focal seizures. It may cause dizziness, double vision, poor coordination, or changes in the heart's electrical rhythm.

\medterm{Lacquer Poisoning} Harm from swallowing or breathing lacquer or its solvents. It may irritate the mouth and lungs, cause dizziness or drowsiness, and cause severe lung injury if liquid is inhaled; contact poison control after significant exposure.

\medterm{Lacrimal} Relating to tears or the lacrimal glands and drainage passages.
\medterm{Lacrimal Apparatus} The glands and passages that make, spread, and drain tears. It includes the tear gland, small openings near the inner eyelids, the tear sac, and the duct that drains tears into the nose.

\medterm{Lacrimal Disease} A disorder of the tear-producing or tear-draining system, including dry eye, dacryocystitis, obstruction, inflammation, or abnormal tearing.

\medterm{Lacrimal Duct} A drainage channel that carries tears from the lacrimal sac into the nasal cavity. Blockage can cause
watering of the eye, recurrent infection, swelling near the inner corner, or discharge.

\medterm{Lacrimal Gland} One of two glands above the outer part of the eye that make the watery part of tears. The tears spread across the eye and then drain through passages near the inner corner.

\medterm{Lacrimal Gland Tumor} A benign or malignant growth arising in the tear-producing lacrimal gland, often causing swelling, pain, or displacement of the eye.

\medterm{Lacrimation} Production and release of tears by the lacrimal glands; excessive tearing may result from irritation, allergy, infection, or blocked tear drainage.

\medterm{Lactam} A ring-shaped chemical structure found in some medicines. Beta-lactam antibiotics use this structure to weaken bacterial cell walls, but allergies and bacterial resistance can limit their use.

\medterm{Lactase} An enzyme in the small intestine that breaks lactose, the natural sugar in milk, into absorbable sugars. Low lactase can cause bloating, cramps, gas, diarrhea, or discomfort after dairy foods.

\medterm{Lactase Deficiency} Alternate index label; see \textit{Lactose intolerance}.

\medterm{Lactate} A product of cellular metabolism measured in blood. An increased level may occur with strenuous exercise, impaired oxygen delivery, infection, medicines, liver disease, or other conditions and must be interpreted clinically.

\medterm{Lactate Dehydrogenase} An enzyme found in many body tissues. Its blood level may rise when cells are damaged, but the test does not identify the cause by itself.

\medterm{Lactate Dehydrogenase Test} A blood test measuring lactate dehydrogenase, an enzyme released when cells are injured; increased activity is nonspecific and may occur with hemolysis, liver disease, muscle injury, infection, or malignancy.

\medterm{Lactation} The production and release of breast milk after childbirth. Hormones stimulate the breasts to make milk and release it during feeding or expression.

\medterm{Lactation (Breastfeeding)} The production and release of breast milk after childbirth. Early milk, called colostrum, is followed by mature milk; hormones help make and release the milk.

\synonyms
Breastfeeding

\medterm{Lactation Disorder} A problem with making, releasing, or transferring breast milk after childbirth. Examples include low supply, excess milk, painful engorgement, mastitis, or difficulty with the baby’s latch; assessment considers both parent and infant.

\medterm{Lactation Mastitis} Alternate index label; see \textit{Mastitis}.

\medterm{Lactation Nutrition} Eating and drinking enough during breastfeeding to support the parent’s health and milk production. Most people need a varied diet and regular fluids; supplements may be needed for a diagnosed deficiency or specific dietary pattern.

\medterm{Lactational Amenorrhea} The temporary absence of menstrual periods during breastfeeding because ovulation is suppressed. Breastfeeding prevents pregnancy reliably only when specific conditions are met, so other contraception may be needed.

\medterm{Lacteals} Tiny lymph vessels in the lining of the small intestine that absorb most digested fats and fat-soluble vitamins. They carry these nutrients through the lymph system before they enter the bloodstream.

\medterm{Lactic Acid} An organic acid produced during normal metabolism, especially when pyruvate is converted to lactate. Blood lactate may rise for many reasons, including impaired oxygen delivery, strenuous exercise, or medicines.

\medterm{Lactic Acid Test} A blood test measuring lactate, which accumulates when production exceeds clearance; it helps assess tissue hypoperfusion, sepsis, severe exercise, metabolic disease, or other causes of lactic acidosis.

\medterm{Lactic Acidosis} Metabolic acidosis associated with increased lactate production or reduced lactate clearance. Causes include impaired tissue oxygen delivery, medicines, liver disease, seizures, and other disorders.

\medterm{Lactobacillus} A genus of bacteria that produces lactic acid and normally forms part of the healthy vaginal and intestinal microbiota; some species are used as probiotics.

\medterm{Lactoferrin} Lactoferrin is an iron-binding glycoprotein in milk and mucosal secretions that limits microbial iron access and contributes to innate immune defense.

\medterm{Lactose} The principal sugar in milk, composed of glucose and galactose. Lactase in the small intestine digests it.

\medterm{Lactose Fermentation} The ability of bacteria to break down lactose and produce acid or gas. It is used in the laboratory to help identify different types of bacteria.

\medterm{Lactose Intolerance} A condition in which the body cannot fully digest lactose, the sugar in milk and dairy products, because it does not make enough lactase. It may cause bloating, gas, abdominal pain, nausea, or diarrhea after consuming lactose.

\synonyms
Lactase Deficiency
Milk Intolerance
Milk Intolerance (Lactose Intolerance)

\medterm{Lactose Tolerance Test} A test for difficulty digesting lactose, the sugar in milk. After lactose is consumed, blood glucose or breath hydrogen is measured; symptoms and results help assess lactose malabsorption but do not diagnose every cause of digestive discomfort.

\medterm{Lactose Tolerance Tests} Tests that evaluate digestion and absorption of lactose by measuring blood glucose response after lactose ingestion or detecting hydrogen in exhaled breath; they may support evaluation of lactose malabsorption.

\medterm{Lactulose} A medicine that softens stool by drawing water into the bowel. It is also used to lower harmful ammonia levels in some people with severe liver disease. Gas and diarrhea can occur.

\medterm{Lacuna} A small space or cavity in tissue, such as a space containing an osteocyte in bone.
\medterm{Lacunar Stroke} A small ischemic stroke caused by blockage of a tiny artery deep in the brain, often producing weakness, numbness, clumsiness, or speech problems.

\medterm{Ladd Bands} Abnormal bands of tissue present at birth that can cross and squeeze the first part of the small intestine when the bowel is in an abnormal position. They may cause green vomiting and blockage in an infant and are released during the Ladd procedure.

\synonyms
Ladd's Peritoneal Bands; Ladd Fibrous Bands

\medterm{Ladd Procedure} An operation for a child whose intestines developed in an abnormal position and may twist. The surgeon untwists the intestine, releases bands causing blockage, and broadens its attachment to reduce future twisting.

\synonyms
Ladd Operation; Ladd's Procedure

\medterm{Laetrile} A semisynthetic amygdalin-related compound promoted as an alternative cancer treatment; credible evidence has not established anticancer benefit, and cyanide poisoning is a serious risk.

\medterm{Lafora Disease} A rare inherited epilepsy that usually begins during adolescence and progressively worsens. Seizures, brief muscle jerks, visual symptoms, memory problems, and loss of physical or thinking abilities may occur.

\medterm{Lagophthalmos} Inability to close the eyelids completely. Exposed cornea can become dry, irritated, scratched, infected, or damaged, especially during sleep, so lubricating treatment and protection may be needed.

\medterm{Lambda} A landmark on the back of the skull where two skull seams meet. It is used to describe head shape and the position of nearby structures, especially in newborn examinations and imaging.
\medterm{Lambert-Eaton Myasthenic Syndrome} A rare condition in which the immune system interferes with nerve signals to muscles. It causes weakness, especially in the hips and shoulders, reduced reflexes, and sometimes dry mouth. It can be associated with small-cell lung cancer.

\medterm{Lamellar} Made of thin layers or plates. The word may describe the organized layers of mature bone, the skin barrier, or structures in the eye and nervous system. Its exact meaning depends on the body structure being described.
\medterm{Lamellar Body} A small storage structure inside certain cells that contains layered fats. In lung cells, it stores and releases surfactant, the substance that helps keep the tiny air sacs open.

\medterm{Lamellar Bone} Mature bone made of organized layers of mineralized tissue. It forms most adult compact and spongy bone and provides strength while being continually remodeled.

\medterm{Lamellar Ichthyosis} An inherited skin disorder present from birth that causes dry skin covered with large, dark, plate-like
scales over much of the body, including the arms and legs. It may reduce sweating and cause overheating,
and can affect the hair, nails, eyes, or palms and soles.

\medterm{Lamin} A structural protein forming part of the nuclear lamina, which supports the inner surface of the cell nucleus.

\medterm{Lamina} A flat piece of bone that forms part of the back of a vertebra. The two laminae help complete the bony arch around the spinal cord and can be removed during some spine operations.

\medterm{Lamina Propria} A thin layer of supporting tissue beneath the inner lining of organs such as the intestine, airways, and bladder. It contains small blood vessels, immune cells, glands, and connective tissue.

\medterm{Laminectomy} An operation that removes part of a spinal bone to make more room for the spinal cord or nerves. It may be used for spinal narrowing, a tumor, or selected disk problems.

\medterm{Laminin} A protein in the supporting framework around cells that helps cells attach to underlying tissue. It contributes to tissue structure, development, repair, and the filtering barrier in some organs.

\medterm{Lamivudine} Lamivudine is a nucleoside reverse-transcriptase inhibitor used with other antiretrovirals for HIV and as an antiviral for chronic hepatitis B; resistance and rare lactic acidosis or hepatic complications require monitoring.

\medterm{Lamotrigine} Lamotrigine is an antiseizure medicine that inhibits voltage-gated sodium channels and is also used for bipolar-maintenance treatment; rapid dose escalation increases the risk of serious rash including Stevens--Johnson syndrome.

\medterm{Lancet} A small, sterile, usually single-use device with a sharp point used to puncture the skin and obtain a capillary-blood sample. Lancets are commonly used for fingertip glucose testing and other point-of-care blood tests; used lancets must be discarded in an appropriate sharps container.

\synonyms
Blood lancet
Fingerstick lancet

\medterm{Landolt Ring} A ring-shaped symbol with a gap used to test visual sharpness. The person identifies which direction the gap faces; smaller symbols indicate better visual acuity.

\medterm{Langer Lines} Natural lines of skin tension that reflect the predominant orientation of collagen fibers in the dermis and may influence surgical scar direction.

\medterm{Langerhans Cell Histiocytosis} A rare disorder in which abnormal immune cells build up in one or more organs, commonly the bones, skin, lungs, lymph nodes, or pituitary gland. Symptoms and treatment depend on which organs are affected.
\medterm{Langerhans Cells} Immune cells in the outer skin that capture foreign substances and help start immune responses. They connect skin defenses with nearby lymph tissue and can contribute to allergic or inflammatory reactions.

\medterm{Langerhans' Cell} An immune cell in the skin that captures foreign substances and helps activate T-cell responses.

\medterm{Langerhans-Cell Histiocytosis} A rare disorder in which abnormal immune cells collect in one or more organs, especially bone, skin, lungs, lymph nodes, or the pituitary gland. Symptoms depend on the organs involved and may include bone pain, rash, cough, or hormone problems.

\medterm{Language Barrier} Difficulty communicating because a patient and healthcare professional do not share a language well enough for safe understanding. A trained interpreter helps explain symptoms, consent, instructions, and risks.

\medterm{Language Development} The gradual development of understanding, sounds, words, grammar, and communication during childhood. Delays may be related to hearing loss, developmental conditions, neurologic disease, or limited language exposure.
\medterm{Language Disorder} A developmental or acquired impairment in understanding, producing, or using spoken, written, or signed language that affects communication and daily function.


\medterm{Language Therapy} Treatment that improves understanding, expression, reading, writing, or other language skills after
developmental or acquired impairment.

\medterm{Lanolin} Lanolin is a waxy substance derived from sheep wool used as an emollient and in topical products; allergy or irritation can occur.

\medterm{Lanolin Poisoning} Swallowing a small amount of lanolin, a wool-derived moisturizer, usually causes little more than stomach upset. Large amounts or products with added medicines can be more dangerous; contact poison control for symptoms or uncertain exposure.

\medterm{Lansoprazole} A medicine that strongly reduces stomach acid. It is used for heartburn, ulcers, and some H. pylori infections; long-term use may affect mineral or vitamin levels and infection risk.

\medterm{Lanthanum} A chemical element whose carbonate salt can bind dietary phosphate and lower phosphate absorption in people with advanced kidney disease.

\medterm{Lanugo} Fine, soft hair that grows all over the body of a fetus and is typically shed before
birth.

\medterm{Laparoscope} A slender instrument with a light and camera used to view the abdominal or pelvic cavity through a small incision during minimally invasive surgery.

\medterm{Laparoscopic Cholecystectomy} Surgical removal of the gallbladder through small abdominal incisions using a laparoscope and instruments, commonly performed for symptomatic gallstones or gallbladder inflammation.

\medterm{Laparoscopic Colectomy} Laparoscopic colectomy is removal of part or all of the colon through small incisions using a camera and specialized instruments.

\medterm{Laparoscopic Gallbladder Removal} Surgery to remove the gallbladder through several small abdominal openings using a camera and instruments. It commonly treats symptomatic gallstones or gallbladder inflammation; bleeding, infection, bile leakage, and injury to bile ducts are possible.

\medterm{Laparoscopic Gastric Banding} A bariatric operation that places an adjustable band around the upper stomach to create a small pouch and restrict food passage; it is less commonly performed than other bariatric procedures.

\medterm{Laparoscopic Hysterectomy} Removal of the uterus using laparoscopic instruments through small abdominal incisions. The cervix, fallopian tubes, and ovaries may or may not also be removed; the approach and extent depend on the indication, anatomy, and surgical plan.



\medterm{Laparoscopy} A type of minimally invasive surgery in which a thin camera is placed through a small cut in the abdomen. Other small cuts may be used to insert surgical instruments.

\medterm{Laparotomy} An operation that uses a larger cut in the abdomen to reach the abdominal or pelvic organs.

\medterm{Lapatinib} Lapatinib is a targeted cancer medicine that blocks HER2 and epidermal-growth-factor receptors and is used for selected advanced breast cancers.

\medterm{Lapatinib Ditosylate} Lapatinib ditosylate is the salt form of lapatinib, a targeted medicine used for selected HER2-positive breast cancers.

\medterm{Large Airway Obstruction} A blockage in the windpipe or large breathing tubes. It can cause noisy breathing, shortness of breath, and difficulty clearing mucus and may be caused by narrowing, a tumor, or weakness of the airway wall.

\medterm{Large Bowel Obstruction} A blockage in the colon, commonly caused by cancer, twisting of the bowel, scar tissue, or diverticular disease. It can cause increasing swelling, cramping, vomiting, and inability to pass stool or gas and may require urgent treatment.

\medterm{Large Bowel Resection} Surgical removal of a segment of the colon or rectum, followed by reconnection of the bowel or creation of a stoma when necessary to treat cancer, obstruction, ischemia, inflammation, or other disease.


\medterm{Large Colon} The large colon is the main portion of the large intestine that absorbs water, stores fecal material, and moves contents toward the rectum.

\medterm{Large For Gestational Age} A fetus or newborn larger than expected for the number of weeks of pregnancy, usually above the 90th percentile. Maternal diabetes, family size, and excess pregnancy weight gain can contribute and may increase birth complications.

\medterm{Large Intestine} The final part of the digestive tract, extending from the small intestine to the rectum. It absorbs water and helps form and move stool.

\medterm{Large Solitary Luteinized Follicle Cyst Of Pregnancy} A usually harmless ovarian cyst that develops during pregnancy as hormone-sensitive ovarian tissue enlarges. It may be found by chance or as a pelvic mass and often settles after pregnancy.


\medterm{Laron Syndrome} A rare inherited growth disorder in which the body cannot respond properly to growth hormone. It causes severe short stature, characteristic facial features, and very low levels of the growth signal needed for normal bone development.


\medterm{Larva Currens From Strongyloides} A rapidly moving, itchy skin rash caused by Strongyloides larvae moving through the skin. It often appears around the buttocks or trunk and may indicate ongoing intestinal infection requiring antiparasitic treatment.

\medterm{Larva Migrans} A condition caused by migrating larval parasites in human tissues, including cutaneous larva migrans in the skin and visceral larva migrans in internal organs.

\medterm{Laryngeal} Relating to the larynx, the organ in the throat containing the vocal cords and helping protect the airway.

\medterm{Laryngeal Cancer} Cancer of the larynx, or voice box. Tobacco, alcohol, and some HPV infections increase risk. Persistent hoarseness, painful or difficult swallowing, noisy breathing, or a neck lump needs prompt examination.

\medterm{Laryngeal Candidiasis} A yeast infection of the voice box, usually causing hoarseness, throat pain, or difficulty swallowing. It is more likely with weakened immunity or inhaled steroid use and may require examination of the larynx and antifungal treatment.

\medterm{Laryngeal Carcinoma} Laryngeal carcinoma is cancer arising in the voice box and may cause hoarseness, swallowing difficulty, pain, or breathing problems.

\medterm{Laryngeal Cartilage} Laryngeal cartilage forms the structural framework of the voice box, maintains the airway, and supports vocal-fold movement.

\medterm{Laryngeal Cleft} A laryngeal cleft is a congenital opening between the airway and esophagus that can allow aspiration and cause feeding or breathing problems.

\medterm{Laryngeal Edema} Swelling of the tissues of the larynx that can narrow the upper airway and cause hoarseness, stridor, or respiratory distress. It may result from allergy, infection, trauma, or airway instrumentation.

\medterm{Laryngeal Fistula} An abnormal passage between the voice box and nearby tissue or the swallowing tract. It can allow food or saliva to enter the airway and cause coughing, repeated chest infections, or breathing problems.

\medterm{Laryngeal Hemorrhage} Bleeding within or around the larynx, which may result from trauma, instrumentation, vocal strain, vascular lesions, infection, or tumor and can threaten the airway.

\medterm{Laryngeal Masks} Laryngeal masks are supraglottic airway devices placed above the vocal cords to maintain ventilation during anesthesia or emergencies.

\medterm{Laryngeal Mucositis} Inflammation and open sores in the lining of the voice box, often after radiation, chemotherapy, infection, or irritation. It can cause hoarseness, throat pain, painful swallowing, or breathing difficulty.

\medterm{Laryngeal Nerve} A branch of the vagus nerve supplying the larynx; injury can cause hoarseness, swallowing difficulty, or breathing problems.

\medterm{Laryngeal Nerve Damage} Injury to the recurrent or superior laryngeal nerves causing hoarseness, vocal-cord
weakness, swallowing difficulty, or airway compromise. Common causes include thyroid or
neck surgery, malignancy, trauma, and compression along the nerve pathway.

\medterm{Laryngeal Obstruction} Partial or complete blockage of the laryngeal airway caused by edema, foreign body, tumor, trauma, infection, vocal-cord dysfunction, or other structural disease.

\medterm{Laryngeal Prosthesis} A device used to help restore speech after removal of the larynx, commonly through a one-way valve placed between the trachea and esophagus.

\medterm{Laryngeal Stridor} A harsh, high-pitched noise caused by turbulent airflow through a narrowed upper airway, often most noticeable during inspiration.

\medterm{Laryngeal Warts (Recurrent Respiratory Papillomatosis)} Benign papillomatous growths caused by HPV in the larynx or respiratory tract. They may cause hoarseness, stridor, or airway obstruction.

\synonyms
Recurrent Respiratory Papillomatosis

\medterm{Laryngeal Web} A laryngeal web is a thin membrane across part of the airway, usually congenital or acquired after injury, causing hoarseness or airway narrowing.

\medterm{Laryngectomee} A person who has undergone laryngectomy, surgical removal of all or part of the larynx; total laryngectomy separates the airway from the mouth and nose.

\medterm{Laryngectomy} Surgery to remove part or all of the larynx, usually for cancer or severe structural disease. After total removal, breathing occurs permanently through an opening in the neck; partial surgery may preserve some voice and swallowing.

\medterm{Laryngectomy Tube} A tube placed in the neck opening after larynx surgery to help keep the airway open and allow suctioning or ventilation when needed. Mucus blockage, bleeding, infection, and tissue overgrowth can occur.

\medterm{Laryngitis} Inflammation of the voice box that causes hoarseness or temporary voice loss. Common causes include viral infection, voice overuse, smoke or other irritation, and reflux.

\medterm{Laryngocele} An air-filled pouch formed when a small space beside the vocal cords becomes enlarged. It may cause hoarseness, a neck swelling that changes with pressure, coughing, or blockage of the airway.

\medterm{Laryngomalacia} A birth condition in which soft tissue above an infant’s vocal cords partly collapses during breathing in, causing noisy breathing. It often improves with growth, but feeding trouble, pauses in breathing, or poor growth needs assessment.

\medterm{Laryngopharyngeal Reflux} Backflow of stomach contents into the throat and voice box. It may cause hoarseness, frequent throat clearing, cough, a lump-like feeling, excess mucus, or throat tightness, often without typical heartburn.

\medterm{Laryngopharynx} The lower part of the throat behind the voice box, leading to the esophagus. It helps direct swallowed food away from the airway and is a site where some throat cancers can develop.
\medterm{Laryngoscope} An instrument used to see the throat and vocal cords, especially while placing a breathing tube. Common blades are curved Macintosh and straight Miller types; use can injure teeth, lips, or throat.

\medterm{Laryngoscopy} An examination of the voice box and nearby throat using a mirror or a thin camera passed through the mouth or nose. It helps investigate hoarseness, swallowing problems, airway narrowing, bleeding, or suspected tumors.

\medterm{Laryngoscopy And Nasolaryngoscopy} An examination of the voice box and nearby airway using a small rigid or flexible camera passed through the mouth or nose. It helps investigate hoarseness, swallowing difficulty, noisy breathing, bleeding, or a suspected growth.

\medterm{Laryngospasm} Sudden tightening of the vocal cords that partly or completely blocks breathing. It can be triggered by saliva, blood, reflux, airway procedures, or anesthesia and may cause noisy breathing, choking, or inability to breathe.

\medterm{Laryngostenosis} Narrowing of the voice box, or larynx, that can restrict breathing. Causes include scarring after intubation, injury, inflammation, tumors, and congenital narrowing; noisy breathing or breathlessness needs medical assessment.

\medterm{Laryngotracheobronchitis (Croup)} Alternate index label for Croup.

\medterm{Larynx} The voice box in the front of the neck between the throat and windpipe. It contains the vocal cords, produces sound, helps keep food and liquid out of the lungs, and regulates airflow during breathing.

\medterm{Larynx Cancer} Laryngeal cancer is a malignant tumor of the voice box and may cause persistent hoarseness, swallowing pain, neck mass, ear pain, cough, or breathing difficulty. Diagnosis requires examination and biopsy; treatment depends on site and stage.

\synonyms
Throat Cancer (Larynx Cancer)

\medterm{Larynx Disease} Any disorder affecting the larynx, including inflammation, infection, structural lesions, vocal-cord dysfunction, neurologic disease, or cancer.

\medterm{Larynx Muscle} A muscle of the voice box that moves the vocal folds and helps protect the airway and produce sound.

\medterm{Larynx Neoplasm} An abnormal growth in the voice box. It may be noncancerous or cancerous and can cause hoarseness, swallowing difficulty, pain, breathing problems, or a neck lump.

\medterm{Laser} A device that produces a concentrated beam of light. Depending on its wavelength and power, it can measure tissue, cut or seal tissue, remove abnormal growths, or treat selected eye and skin conditions.
\medterm{Laser (Vascular)} A laser system used to treat selected blood-vessel problems or tissue blocking a vessel. The energy is delivered through the skin or a catheter and must be carefully targeted to avoid damage to nearby tissue.
\synonyms
Vascular laser

\medterm{Laser Ablation} The use of focused laser energy to heat, destroy, or seal targeted tissue. The technique is used in selected vascular, ophthalmic, dermatologic, and other procedures.

\medterm{Laser Assisted Uvula Palatoplasty} A procedure using laser energy to reshape or remove selected soft-palate and uvula tissue, intended mainly to reduce snoring in carefully selected patients.

\synonyms
LAUP


\medterm{Laser Fiber} A thin optical fiber that carries laser energy into the body during treatment. It can break urinary stones, destroy selected tumors, seal blood vessels, or treat other targeted tissues.

\medterm{Laser Photocoagulation} A laser treatment that seals or destroys leaking or abnormal blood vessels, especially in the retina, helping limit bleeding, swelling, or damaging new vessel growth.

\medterm{Laser Photocoagulation - Eye} An eye treatment that uses laser heat to seal leaking blood vessels or destroy abnormal retinal tissue, helping treat selected retinal diseases.
\medterm{Laser Prostatectomy} A procedure using laser energy to remove or vaporize enlarged prostate tissue and improve urine flow. It is performed through the urethra and may cause bleeding, infection, temporary urgency, or other urinary problems.

\medterm{Laser Surgery} Laser surgery uses a focused beam of light to cut, remove, seal, or reshape tissue in selected procedures.

\medterm{Laser Surgery For The Skin} Use of concentrated light to remove, vaporize, seal, or reshape selected skin tissue. It may treat scars, abnormal blood vessels, unwanted hair, or some lesions; redness, pigment change, burns, scarring, and infection are possible.

\medterm{Laser Therapy} Laser therapy uses controlled light energy to treat selected skin, eye, vascular, dental, or other medical conditions.

\medterm{Laser Therapy For Cancer} Use of an intense, focused beam of light to destroy or shrink selected cancer tissue or control bleeding and blockage. It is useful only for certain tumors and sites; treatment may be combined with surgery, medicines, or radiation.

\medterm{Laser Trabeculoplasty} A laser treatment for some forms of open-angle glaucoma that improves drainage of fluid from the eye and lowers eye pressure. The effect may be temporary, and continued examinations or other treatment may still be needed.

\medterm{Laser-Assisted In Situ Keratomileusis} A refractive surgical procedure that reshapes the cornea using a laser to correct short-sightedness, long-sightedness, or astigmatism, thereby reducing dependence on glasses or
contact lenses.

\synonyms
LASIK

\medterm{LASIK} A laser operation that reshapes the clear front surface of the eye to correct short-sightedness, long-sightedness, or astigmatism. It can reduce dependence on glasses, but dry eye, glare, infection, under-correction, and other vision problems are possible.

\medterm{LASIK Eye Surgery} A laser operation that reshapes the cornea to correct short-sightedness, long-sightedness, or astigmatism and reduce dependence on glasses. Dry eyes, glare, infection, under- or over-correction, and rarely serious vision loss can occur.


\medterm{Lassa Fever} A viral illness that can cause fever, weakness, vomiting, bleeding, and damage to several organs. It is usually spread by contact with urine or droppings of infected rodents in parts of West Africa and sometimes by contact with body fluids.

\medterm{Lassa Virus} An arenavirus that causes Lassa fever, a potentially severe hemorrhagic illness acquired mainly through contact with infected rodent excreta or body fluids.

\medterm{Lassitude} A state of marked fatigue, weakness, or lack of energy that may occur with systemic illness, anemia, endocrine disease, infection, medication effects, or psychological conditions.

\medterm{Last Menstrual Period} The first day of the most recent normal menstrual bleeding. Clinicians use it to estimate how many weeks a pregnancy has lasted and the expected delivery date, but the estimate is less reliable with irregular cycles or uncertain dates.

\medterm{Latanoprost} Latanoprost is a prostaglandin F2-alpha analogue eye drop that increases uveoscleral outflow and lowers intraocular pressure in glaucoma or ocular hypertension; iris pigmentation and eyelash growth may occur.

\medterm{Latarjet Procedure} Shoulder-stabilization surgery that moves part of the coracoid bone, with attached tissue, to the front edge of the shoulder socket. It helps prevent repeated dislocation when bone loss or failed previous treatment makes other repairs unsuitable.

\medterm{Late Decelerations} A repeated, gradual slowing of a baby’s heart rate that starts after a uterine contraction begins and recovers after the contraction ends. It can indicate reduced oxygen transfer through the placenta and requires assessment during labor.

\medterm{Late Effects} Health problems appearing months or years after cancer treatment, including heart, lung, hormone, fertility, nerve, bone, or second-cancer complications.

\medterm{Late-Onset Hypogonadism} A condition in which an older man has symptoms of low testosterone and repeatedly low early-morning blood levels. Symptoms may include reduced sexual desire, erection problems, tiredness, reduced muscle strength, increased body fat, weak bones, or low mood. Other illnesses and medicines should also be assessed before treatment.

\synonyms
LOH; Age-Related Testosterone Deficiency; Testosterone Deficiency Syndrome

\medterm{Late Gadolinium Enhancement} A heart-MRI finding in which contrast dye remains longer in areas of heart muscle that are scarred, injured, or inflamed. Its pattern helps doctors estimate the cause and extent of heart damage.
\synonyms
LGE; Delayed Contrast Enhancement; Myocardial Scar Imaging

\medterm{Late Preterm (34-36 Weeks)} Births occurring from 34 weeks 0 days through 36 weeks 6 days of gestation. These infants have
greater risks than term infants of respiratory, temperature, glucose, jaundice, feeding,
and readmission problems.

\synonyms
34-36 weeks

\medterm{Latent} Present but inactive, hidden, or not currently producing symptoms, as in latent tuberculosis infection or a latent virus.

\medterm{Latent Autoimmune Diabetes In Adults} Diabetes caused by the immune system gradually damaging insulin-producing cells in an adult. It may first resemble type 2 diabetes, but insulin production usually declines over time and insulin treatment may eventually be needed.

\synonyms
Type 1.5 diabetes; Late-onset autoimmune diabetes of adulthood

\medterm{Latent Infection} Persistent infection in which an infectious agent remains in a host in a nonreplicating
or intermittently replicating state without active clinical disease. Reactivation can
occur when immune control is impaired, as in latent tuberculosis.

\medterm{Latent Phase Of Labor} The early part of the first stage of labor, beginning with regular contractions and gradual
cervical change before active labor. The transition to active labor is assessed clinically,
and the duration varies widely between individuals.

\medterm{Latent Syphilis} Latent syphilis is syphilis infection without current symptoms, identified by testing and classified by duration to guide treatment and follow-up.

\medterm{Latent Tuberculosis} Infection with Mycobacterium tuberculosis in which the bacteria remain inactive and cause no symptoms but may later become active.

\medterm{Latent Tuberculosis Infection} An immune response to Mycobacterium tuberculosis without clinical or microbiologic evidence of active
disease. It is usually asymptomatic and not contagious, but it can reactivate; treatment
decisions depend on age, immune status, exposure risk, test results, and drug interactions.

\medterm{Lateral} Toward or located at the side, away from the body’s midline. For example, the ears are lateral to the nose; the opposite term is medial.

\medterm{Lateral Canthotomy And Cantholysis} An emergency operation that opens the outer corner of the eyelids and releases tight supporting tissue around the eye. It lowers dangerous pressure after severe eye injury or bleeding and may protect vision.

\medterm{Lateral Collateral Ligament Injury} A stretch or tear of the strong ligament on the outer side of the knee, often caused by a blow to the inner knee or a force pushing the lower leg inward. It causes outer-knee pain and may make the knee unstable.

\medterm{Lateral Cuneiform} One of three small bones in the middle of the foot. It lies between the other cuneiform bones and the cuboid, joins the third toe bone, and helps support the foot’s arch.

\medterm{Lateral Decubitus Position} A side-lying position in which the patient rests on the nonoperative side to provide access to the chest, hip, or flank.

\medterm{Lateral Elbow Tendinopathy} Alternate index label; see \textit{Tennis elbow}.

\medterm{Lateral Epicondylitis} Alternate index label; see \textit{Tennis elbow}.

\medterm{Lateral Epicondylitis (Tennis Elbow (Lateral Epicondylitis))} Alternate index label for Tennis Elbow (Lateral Epicondylitis).

\medterm{Lateral Epicondylosis} Alternate index label; see \textit{Tennis elbow}.

\medterm{Lateral Meniscus} The lateral meniscus is a crescent-shaped cartilage pad on the outer side of the knee that distributes load and stabilizes movement.

\medterm{Lateral Rotation} Turning a limb or body part away from the body’s midline. Rotating the hip so the knee and foot point outward is one example; limited or painful movement may indicate injury.

\medterm{Lateral Sclerosis} Damage affecting nerve pathways along the side of the spinal cord that carry movement signals. It may cause increasing stiffness and weakness, as in some motor-neuron diseases; the exact diagnosis requires the associated clinical findings.

\medterm{Lateral Traction} Lateral traction is a pulling force applied sideways to a body part to improve alignment, reduce a dislocation, or relieve pressure.

\medterm{Lateral Ventricle} One of two fluid-filled spaces inside the brain that produce and contain cerebrospinal fluid, helping protect and support the brain.

\medterm{Lateral X-Ray} An X-ray image taken from the side, providing a lateral view of the body part being examined.

\medterm{Latex Agglutination Test} A laboratory test in which tiny coated particles clump together when they meet a matching substance. The reaction can help detect selected germs, antibodies, or proteins in a patient sample.

\medterm{Latex Allergies - For Hospital Patients} Latex allergy can cause contact dermatitis, hives, wheeze, or anaphylaxis; hospitals reduce risk with latex-free supplies and clear documentation of the allergy.

\medterm{Latex Allergy} An immune-mediated hypersensitivity reaction to natural rubber latex. Contact may cause
dermatitis, urticaria, rhinitis, asthma, or anaphylaxis, especially in healthcare and
occupational settings.

\synonyms
Allergy Latex (Latex Allergy)
Allergy, Latex

\medterm{Lathyrism} A neurologic disorder caused by prolonged or excessive ingestion of certain Lathyrus legumes, especially grass pea, characterized primarily by progressive spastic weakness of the legs.

\medterm{Latissimus Dorsi} A broad back muscle that extends and internally rotates the arm and assists movements such as climbing.
\medterm{Latissimus Dorsi Muscle} A broad muscle covering much of the back that pulls the arm down, backward, and toward the body. It is used in climbing, swimming, and rowing and may also be used to reconstruct tissue during surgery.

\medterm{Lavage} Washing a body space, organ, or wound with sterile fluid. It may be used to remove blood, infected material, or harmful substances, or to collect a sample; the method and risks depend on the part of the body treated.

\medterm{Lavage Fluid} Liquid used to wash a body space, wound, or organ during a medical procedure, or collected afterward for laboratory analysis. Its contents can help detect infection, bleeding, cancer cells, inflammation, or abnormal proteins.


\medterm{Lawton IADL Scale} A tool measuring ability to manage complex daily tasks such as telephone use, shopping, cooking, housekeeping, transportation, medicines, and finances.

\medterm{Laxative Overdose} Taking too much laxative can cause severe diarrhea, vomiting, dehydration, and abnormal body salts. Some products can disturb heart rhythms or damage the bowel; urgent medical advice is needed for severe symptoms, children, or large exposures.

\medterm{Laxatives For Constipation} Laxatives treat constipation through fiber and water retention, stool softening, osmotic effects, or increased bowel contractions. Choice depends on cause, duration, age, medicines, kidney function, and obstruction risk; persistent bleeding, vomiting, or severe pain needs evaluation.

\synonyms
Constipation (Laxatives For Constipation)


\medterm{Lazy Eye} Reduced vision caused by abnormal visual development during childhood, usually in one eye but sometimes in both. It may result from eye misalignment, unequal focusing, or an obstruction such as a cataract; the medical term is amblyopia.

\synonyms
amblyopia.

\medterm{Lazy Eye (Amblyopia)} Reduced vision in one or both eyes caused by abnormal visual development during childhood, often related to misalignment, unequal focusing, or visual obstruction.

\medterm{LCIS} Alternate index label; see \textit{Lobular carcinoma in situ}.

\medterm{LDH Isoenzyme Blood Test} A blood test separating lactate dehydrogenase isoenzymes to help suggest the tissue source of an elevated total LDH, although more specific tests are preferred for most conditions.

\medterm{LDL} Low-density lipoprotein, a blood particle that carries cholesterol to tissues. Too much LDL can contribute to fatty buildup in artery walls and increase the risk of heart attack and stroke; risk-based treatment may lower it.

\medterm{LDL Test} A blood test measuring low-density lipoprotein cholesterol, a major atherogenic lipoprotein used with the lipid profile to estimate cardiovascular risk and guide lipid-lowering treatment.

\medterm{Lead} A toxic heavy metal that can harm the nervous system, blood-forming system, kidneys, cardiovascular system, and reproductive system. Exposure may occur through contaminated paint dust or soil, drinking water, workplaces, hobbies, imported products, or other sources; children and fetuses are especially vulnerable to neurodevelopmental effects.

\medterm{Lead - Nutritional Considerations} Adequate iron, calcium, and vitamin C intake may help reduce lead absorption in people who are deficient, especially children, but nutrition does not remove lead from the body or replace exposure control. Suspected lead exposure requires evaluation of the source and appropriate blood lead testing.

\medterm{Lead And Tap Water} Drinking water can become contaminated with lead when it contacts lead-containing service lines, plumbing, fixtures, or solder. Risk reduction may require testing, use of an appropriately certified filter, flushing or other interim measures, and repair or replacement of the source according to local public-health guidance; boiling water does not remove lead.

\medterm{Lead Levels - Blood} A blood lead test measures the amount of lead in blood, usually reported in micrograms per deciliter (µg/dL). Results are interpreted with age, symptoms, exposure history, and testing method; an elevated capillary result may require confirmation with a venous sample. Management focuses on removing the exposure and arranging medical or public-health follow-up, with chelation reserved for selected cases.

\medterm{Lead Poisoning} Harm caused by exposure to lead, which can damage the brain, nerves, blood, kidneys, and reproductive system. Children may develop learning, behavior, growth, or developmental problems. Sources include old paint and dust, contaminated soil or water, workplace or hobby materials, and some imported products. Diagnosis uses blood lead testing; treatment starts by stopping exposure, with chelation for selected severe cases.

\medterm{Lead-Time Bias} The false appearance that screening improves survival because it finds a disease earlier, even though the person dies at the same time as without earlier diagnosis.



\medterm{Learned Insomnia} Ongoing difficulty sleeping maintained by worry, sleep-related habits, and learned associations that make the bed or bedtime trigger alertness instead of sleep.

\medterm{Learning Difficulties} Persistent problems acquiring or using academic skills such as reading, writing, or mathematics, not explained solely by low intelligence or inadequate education.
\medterm{Learning Disabilities} Persistent difficulties acquiring or using skills such as reading, writing, mathematics, language, attention, or coordination despite appropriate opportunity and instruction; assessment considers development, education, and other neurologic or sensory factors.

\medterm{Learning Disorder} A neurodevelopmental disorder causing persistent difficulty acquiring or using academic skills such as reading, writing, or mathematics despite appropriate instruction.


\medterm{Leber Congenital Amaurosis} A rare inherited disorder of the retina that causes severe vision loss from birth or early infancy. Babies may have wandering eye movements, sensitivity to light, or unusual responses of the pupils; genetic testing can help confirm the cause.

\medterm{Leber Hereditary Optic Neuropathy} An inherited disorder of the energy-producing parts of cells that damages the optic nerves. It usually causes painless, rapidly worsening loss of central vision in one eye followed by the other, often in young adults. The altered mitochondrial DNA is passed through the mother; genetic testing can confirm the diagnosis.

\synonyms
LHON

\medterm{Leech} A blood-feeding annelid used in selected reconstructive procedures to relieve venous congestion; leech bites can cause bleeding and infection.
\medterm{Leech Therapy} Carefully controlled use of medicinal leeches to remove congested blood from selected tissue, especially after reconstructive surgery. Their saliva also prevents clotting. Bleeding, infection, anemia, and allergic reactions are possible, so clinical monitoring is essential.

\medterm{Leflunomide} Leflunomide is a disease-modifying antirheumatic drug that inhibits dihydroorotate dehydrogenase and lymphocyte proliferation; it treats rheumatoid or psoriatic arthritis but can cause hepatotoxicity, cytopenias, hypertension, and fetal harm.

\medterm{Left Anterior Descending Artery} A major branch of the left coronary artery that travels in the anterior interventricular
groove and supplies the anterior left ventricle, much of the interventricular septum,
and the apex. Its occlusion can cause extensive anterior myocardial infarction.

\medterm{Left Anterior Descending Coronary Artery} A major heart artery supplying the front wall, septum, and much of the left ventricle; blockage can cause a serious heart attack.

\synonyms
LAD artery

\medterm{Left Atrial Appendage Occlusion} A procedure that seals a small pouch of the heart where blood clots often form in people with atrial fibrillation. A closure device is placed through a catheter, usually from a vein in the groin. It may reduce stroke risk when long-term blood-thinning medicine is unsuitable, but short-term medication and follow-up are still needed.

\synonyms
LAAO; Left Atrial Appendage Closure; Watchman Procedure

\medterm{Left Atrial Myxoma} A usually benign primary cardiac tumor, commonly arising from the interatrial septum in the left
atrium. It may obstruct mitral inflow or cause embolic and constitutional symptoms such as
fever, weight loss, or fatigue.

\medterm{Left Atrium} The upper-left chamber of the heart. It receives oxygen-rich blood from the lungs and passes it through the mitral valve into the left ventricle.

\medterm{Left Bundle Branch} Part of the heart’s electrical wiring that carries signals through the left lower chamber and helps it contract. Damage or blockage can delay the heartbeat signal and appear on an electrocardiogram.

\medterm{Left Circumflex Coronary Artery} A heart artery curving around the left side and back, supplying parts of the left ventricle and sometimes the lower heart wall.

\synonyms
Circumflex artery

\medterm{Left Heart Catheterization} A procedure that passes a thin tube into the left side of the heart or nearby arteries to measure pressures, assess valves, or look for blocked coronary arteries.
\medterm{Left Heart Ventricular Angiography} An imaging study that injects contrast into the left ventricle to assess ventricular contraction, wall motion, ejection fraction, and selected structural abnormalities.

\medterm{Left Hemiparesis, Gaze Deviated To The Right} Weakness on the left side with both eyes looking to the right. This pattern can occur with a serious injury to the right side of the brain, including stroke, but examination and brain imaging are needed to identify the cause.

\medterm{Left Hypochondriac Region} The upper-left area of the abdomen, beneath the lower ribs. It overlies part of the stomach, spleen, pancreas, and colon; pain or swelling there may arise from digestive, splenic, or nearby chest disease.

\medterm{Left Iliac Region} The lower-left area of the abdomen, overlying part of the large intestine and urinary tract. Pain there may occur with diverticular disease, bowel inflammation, urinary problems, or other pelvic conditions.

\medterm{Left Main Coronary Artery} The short vessel arising from the left aortic sinus and dividing principally into the
left anterior descending and circumflex arteries. Significant left-main disease
jeopardizes a large proportion of left-ventricular myocardium.

\medterm{Left Testicular Torsion} Acute rotation of the left testis on its spermatic cord, obstructing blood flow and
causing ischemia. It is a urological emergency presenting with sudden severe left
scrotal pain, nausea, vomiting, and a high-riding, horizontally oriented testis with
absent cremasteric reflex.

\medterm{Left Ventricle} The lower-left chamber of the heart. It receives oxygen-rich blood from the left atrium and pumps it through the aortic valve into the aorta and throughout the body.

\medterm{Left Ventricular Hypertrophy} Thickening of the muscle wall of the heart's main pumping chamber. It is often caused by long-term high blood pressure or a narrowed aortic valve and can make the heart stiff, leading to breathlessness, rhythm problems, or heart failure.

\medterm{Left Ventricular Noncompaction} A rare heart-muscle disorder in which the main pumping chamber has an unusually spongy inner muscle layer. It may cause no symptoms or lead to heart failure, abnormal rhythms, blood clots, or stroke.


\medterm{Left-Sided Superior Vena Cava} A birth-related variation in which a large vein on the left side of the chest drains blood through an unusual route to the heart. It usually causes no symptoms but matters when placing central lines or performing heart surgery.

\medterm{Left-Ventricular Assist Device} An implanted mechanical pump that helps the heart's main pumping chamber send blood around the body. It may support people with severe heart failure while they await a transplant or as long-term treatment; infection, bleeding, and blood clots are important risks.

\medterm{Leg} The lower limb, often meaning the region between the knee and ankle, although it may also refer broadly to the entire lower extremity.


\medterm{Leg CT Scan} A series of detailed X-ray pictures of the leg combined by a computer. It can show complex fractures, bone alignment, infection, tumors, blood-vessel problems, and selected soft-tissue injuries.

\medterm{Leg Fracture} Alternate index label; see \textit{Broken leg}.

\medterm{Leg Lengthening And Shortening} Operations that gradually increase or decrease the length of a leg bone to correct a major difference between the legs or a deformity. Recovery is prolonged and may involve an external or internal support device and rehabilitation.

\medterm{Leg MRI Scan} Magnetic resonance imaging of the leg using a magnetic field and radio waves to evaluate muscles, tendons, ligaments, nerves, marrow, joints, and soft-tissue masses.

\medterm{Leg Or Foot Amputation} Surgical removal of part or all of a leg or foot, usually because of severe injury, infection, poor blood flow, or cancer. Rehabilitation supports mobility afterward.
\medterm{Leg Or Foot Amputation - Dressing Change} Changing and examining the dressing over a leg or foot amputation wound. The dressing protects the wound, absorbs drainage, and allows assessment for bleeding, infection, poor healing, pressure injury, or separation of the wound edges.

\medterm{Leg Pain} Pain in the lower limb caused by muscle or tendon injury, joint disease, nerve problems, vascular disease, infection, or other conditions.

\medterm{Leg Pain With Cramping} Painful tightening or spasms in a leg muscle; causes include exertion, dehydration, medicines, nerve or muscle disorders, and reduced arterial blood flow causing intermittent claudication.

\medterm{Leg Swelling} Enlargement or puffiness of one or both legs caused by fluid retention, venous disease, injury, infection, inflammation, or a blood clot.

\medterm{Leg Ulcer} An open sore on the leg that fails to heal normally, often caused by poor venous circulation, arterial disease,
diabetes, pressure, or inflammation. Pain, spreading redness, fever, or black tissue requires
prompt assessment.

\medterm{Legg-Calve-Perthes Disease} A childhood hip disorder in which the blood supply to the top of the thighbone is temporarily reduced. The bone can weaken and change shape, causing a limp, hip or knee pain, stiffness, and limited movement.

\medterm{Legg-Calvé-Perthes Disease} A childhood disorder caused by temporary loss of blood supply to the femoral head, resulting in avascular necrosis and hip pain or limp.
\medterm{Legg-Calvé-Perthes Disease (Perthes Disease)} A childhood disorder in which the blood supply to the femoral head is temporarily impaired, causing avascular necrosis, hip pain, stiffness, and possible deformity.

\synonyms
Perthes Disease

\medterm{Legg–Calvé–Perthes Disease} A childhood hip disorder in which reduced blood flow temporarily damages the rounded top of the thighbone. It can cause a limp, hip or knee pain, stiffness, and later deformity as the bone heals.

\medterm{Legionella} Legionella is a group of water-associated bacteria that can cause Legionnaires' disease, a severe pneumonia, or Pontiac fever, a milder illness.

\medterm{Legionella Longbeachae} Legionella longbeachae is a Legionella species associated particularly with exposure to contaminated potting mix and capable of causing pneumonia.

\medterm{Legionella Pneumonia} Pneumonia caused by Legionella bacteria, usually after inhaling contaminated water droplets. It can cause high fever, cough, breathlessness, diarrhea, confusion, and low oxygen and requires antibiotic treatment.

\medterm{Legionella Pneumonia (Legionnaire Disease And Pontiac Fever)} Alternate index label for Legionnaire Disease and Pontiac Fever.

\medterm{Legionella Pneumophila} A bacterium found in water systems that can cause severe pneumonia after a person inhales contaminated water droplets. It may also cause a milder fever illness; older age, smoking, and weakened immunity increase risk.

\medterm{Legionellosis} An infection caused by Legionella bacteria, ranging from a mild flu-like illness to Legionnaires’ disease with pneumonia.

\medterm{Legionnaire Disease And Pontiac Fever} Legionnaires disease is a severe pneumonia caused by Legionella, while Pontiac fever is a milder flu-like illness from the same bacteria. Exposure to contaminated water systems is typical; pneumonia, confusion, low oxygen, or high fever requires urgent treatment.

\synonyms
Legionella Pneumonia (Legionnaire Disease And Pontiac Fever)
Legionnaires Disease (Legionnaire Disease And Pontiac Fever)

\medterm{Legionnaires Disease} A severe pneumonia caused by Legionella bacteria, usually acquired by inhaling contaminated water aerosols. It may also cause gastrointestinal symptoms, confusion, and low serum sodium.

\medterm{Legionnaires Disease (Legionnaire Disease And Pontiac Fever)} Alternate index label for Legionnaire Disease and Pontiac Fever.

\medterm{Legionnaires' Disease} A severe form of pneumonia caused by Legionella pneumophila, a Gram-negative aerobic
bacillus that is found naturally in freshwater environments (lakes, streams) and
artificially in water systems (air conditioning cooling towers, hot tubs, decorative
fountains, plumbing systems, showers, humidifiers).

\medterm{Legs Restless (Restless Leg Syndrome)} Alternate index label for Restless Leg Syndrome.

\medterm{Leigh Disease} A progressive mitochondrial disorder of early life characterized by symmetric lesions of the brainstem and basal ganglia, developmental regression, movement abnormalities, and lactic acidosis.

\medterm{Leiomyoma} A usually noncancerous tumor made of smooth muscle. It most often develops in the uterus, where it is called a fibroid, but it can also occur in the skin, digestive tract, or other organs.

\medterm{Leiomyoma (Fibroid)} A benign tumor of uterine smooth muscle that may cause heavy menstrual bleeding, pelvic pressure, pain, or no symptoms.

\synonyms
Fibroid

\medterm{Leiomyomatosis} Leiomyomatosis is development of multiple benign smooth-muscle tumors or proliferations, which may affect the uterus, skin, lungs, or other sites.

\medterm{Leiomyomatosis Peritonealis Disseminata} A rare benign condition with multiple smooth-muscle nodules scattered over the peritoneum, often associated with pregnancy or hormonal stimulation and sometimes mimicking peritoneal carcinomatosis.

\medterm{Leiomyosarcoma} A malignant tumor of smooth muscle, most often arising in the uterus, retroperitoneum,
or gastrointestinal tract.

\medterm{Leiomyosarcoma Of The Vulva} A rare malignant smooth-muscle tumor of the vulva, usually presenting as a solitary enlarging mass. Histologic atypia, mitotic activity, necrosis, and invasion help determine malignant behavior.

\medterm{Leishmania} A group of single-celled parasites transmitted mainly by infected sandflies. Different species cause skin, mucosal, or visceral leishmaniasis, ranging from localized ulcers to life-threatening illness.


\medterm{Leishmania Donovani} A parasite spread by infected sandflies that causes visceral leishmaniasis. It can infect organs such as the spleen, liver, and bone marrow, causing prolonged fever, weight loss, enlarged organs, and reduced blood-cell counts.

\medterm{Leishmania Infection} A parasitic infection caused by \textit{Leishmania} protozoa and transmitted by infected sandflies; it can cause cutaneous, mucocutaneous, or visceral leishmaniasis.

\medterm{Leishmania Major} A parasite commonly causing cutaneous leishmaniasis in parts of Africa, Asia, and the Middle East. Sandfly transmission produces one or more skin ulcers that may heal slowly and leave scars.

\medterm{Leishmania Mexicana} A parasite that causes mainly cutaneous leishmaniasis in the Americas. Infection after a sandfly bite may produce a persistent skin nodule or ulcer and, less commonly, involve mucous membranes.

\medterm{Leishmania Tropica} A sandfly-transmitted parasite that commonly causes localized cutaneous leishmaniasis. It produces chronic skin sores that may heal over months and leave permanent scars.

\medterm{Leishmaniasis} An infection caused by Leishmania parasites transmitted by sandflies. It may produce skin ulcers, damage to the nose or mouth, or a life-threatening illness affecting organs such as the spleen, liver, and bone marrow.

\medterm{Lemaitre'S Sign} Pain or discomfort felt in the urethra or penis during a rectal examination. It has traditionally been associated with inflammation near the prostate or back part of the urinary passage but is not specific for one diagnosis.
\medterm{Lemierre Syndrome} A potentially severe oropharyngeal infection complicated by septic thrombophlebitis of the internal jugular vein and frequently metastatic infection, classically caused by Fusobacterium necrophorum.



\medterm{Lenalidomide} Lenalidomide is an immunomodulatory anticancer drug used for multiple myeloma and selected myelodysplastic or lymphoma disorders; myelosuppression, thrombosis, infection, and severe teratogenicity are important risks.

\medterm{Length-Time Bias} The tendency of screening to preferentially detect slower-growing or less aggressive
disease because it remains detectable for a longer preclinical period. It can exaggerate
the apparent benefit of screening.

\medterm{Lengthening Of Limb} A surgical or orthopedic process that gradually increases the length of a bone, often using controlled distraction after osteotomy.
\medterm{Lennox-Gastaut Syndrome} A severe childhood-onset developmental and epileptic encephalopathy with multiple seizure types, characteristic slow spike-and-wave EEG patterns, cognitive impairment, and often treatment-resistant seizures.

\medterm{Lenograstim} A recombinant granulocyte colony-stimulating factor that increases neutrophil production and is used to prevent or treat some forms of neutropenia.

\medterm{Lens} A clear, flexible structure behind the iris and pupil that bends light onto the retina. Muscles change its shape for near or distant focus; clouding causes cataract, while displacement can blur vision or raise eye pressure.

\medterm{Lens Coloboma} A lens coloboma is a congenital segmental defect or notch in the eye lens caused by incomplete development.

\medterm{Lens Disease} A disorder of the eye lens, including cataract, dislocation, inflammation-related opacity, or developmental abnormality, that may impair focusing or vision.

\medterm{Lens Dislocation} Displacement of the crystalline lens from its normal position because of trauma, zonular weakness, surgery, or connective-tissue disease; it can cause blurred vision, glaucoma, or inflammation.

\medterm{Lens Prescription} Written instructions specifying the lens power, correction, and other features needed for eyeglasses or contact lenses in each eye.
\medterm{Lens Subluxation} Partial displacement of the crystalline lens while some zonular attachments remain, potentially causing blurred or fluctuating vision, refractive change, or glaucoma.

\medterm{Lensometer} An instrument measuring the prescription, power, and optical alignment of eyeglass lenses. It helps verify new glasses, assess existing lenses, and check whether corrective strength matches the prescribed values.

\synonyms
Lensmeter; focimeter


\medterm{Lentiform Nucleus} A group of deep brain structures involved in controlling movement. It includes the putamen and globus pallidus and works with other brain regions to help start and regulate actions.

\medterm{Lentigines, Solar} Alternate index label; see \textit{Age spots (liver spots)}.

\medterm{Lentigo} A lentigo is a sharply defined flat brown skin spot caused by increased melanocytes or pigment and may be benign or, rarely, atypical.

\medterm{Lentigo Maligna} A melanoma in situ that occurs on chronically sun-damaged skin, most commonly on the
face of elderly individuals. It presents as a large, irregularly pigmented, slow-growing
macule with variegated shades of brown and black. The radial growth phase may persist
for years before invasion occurs.

\medterm{Lentigo Maligna Melanoma} An invasive melanoma arising in a lentigo maligna on chronically sun-damaged skin, usually of the head or neck; it presents as an enlarging irregular pigmented patch or plaque.


\medterm{LEOPARD Syndrome} A rare inherited disorder that can cause many dark skin spots, abnormal heart rhythms, widely spaced eyes, narrowing of the pulmonary valve, genital differences, slow growth, and hearing loss. The combination and severity vary between people.

\medterm{Leopold Maneuvers} A systematic abdominal palpation technique used to determine fetal position and
presentation. Four maneuvers: (1) fundal grip identifies fetal lie (head vs breech), (2)
umbilical grip locates the fetal back, (3) Pawlik grip identifies the presenting part
and its position, (4) pelvic grip confirms presentation and attitude.

\medterm{Lepidoptera} The insect order containing butterflies and moths; contact with some caterpillars can cause skin, eye, or systemic reactions.

\medterm{Lepirudin} Lepirudin is a direct thrombin inhibitor formerly used as an anticoagulant in heparin-induced thrombocytopenia; bleeding and kidney clearance are important considerations.

\medterm{Lepromatous Leprosy} A widespread form of leprosy caused by a large number of bacteria. It can produce many skin patches, nodules, thickened skin, loss of feeling, nerve damage, and facial changes, with a high risk of spread to other tissues if untreated.

\medterm{Lepromin Skin Test} A skin test using harmless material from the leprosy bacterium to show how strongly the immune system reacts. It may help classify the type of leprosy but cannot by itself diagnose active infection.

\medterm{Leprosy (Hansen Disease)} A long-term infection caused mainly by bacteria that affect the skin and cooler nerves in the arms, legs, and face. It can cause pale or reddish patches with reduced feeling, thickened nerves, weakness, and eye or nose problems. A skin sample or other tests may confirm it. Treatment uses several antibiotics, commonly rifampicin, dapsone, and clofazimine, to prevent lasting nerve damage and disability.

\medterm{Leptin} A hormone produced mainly by adipose tissue that signals stored-energy status to the brain and influences appetite and energy balance.

\medterm{Leptin Receptor} A cell-surface protein that receives signals from leptin, a hormone made mainly by fat tissue. It helps the brain regulate appetite, energy use, growth, and reproductive function in response to stored body energy.

\medterm{Leptomeningeal} Relating to the pia mater and arachnoid mater, the two thin inner layers covering the brain and spinal cord; cancer or infection can spread in these layers.

\medterm{Leptomeningeal Metastasis} Spread of cancer to the thin membranes and fluid surrounding the brain and spinal cord. It can cause headache, vomiting, weakness, confusion, seizures, or changes in vision and requires specialist treatment.

\medterm{Leptomeningeal Sarcoma} Leptomeningeal sarcoma is a rare malignant tumor involving the membranes and fluid spaces around the brain or spinal cord.

\medterm{Leptomeninges} The arachnoid mater and pia mater, the two delicate inner membranes covering the brain and spinal cord.

\medterm{Leptomeningitis} Inflammation or infection of the leptomeninges, the arachnoid and pia mater covering the brain and spinal cord.

\medterm{Leptospira} A genus of spiral-shaped bacteria that includes the causes of leptospirosis, usually acquired through water or soil contaminated with infected animal urine.

\medterm{Leptospirosis} An infection caused by Leptospira bacteria, usually acquired through water or soil contaminated with infected animal urine. It may cause fever, headache, muscle aches, red eyes, vomiting, or severe kidney and liver failure.

\synonyms
Weils Disease (Leptospirosis)

\medterm{Leriche Syndrome} Aortoiliac occlusive disease causing reduced blood flow to the pelvis and lower limbs.
The classic presentation includes buttock or thigh claudication, diminished femoral
pulses, and erectile dysfunction, although symptoms vary with the extent of collateral
circulation.



\medterm{Lesch-Nyhan Syndrome} An inherited disorder of purine metabolism caused by severe HPRT enzyme deficiency, leading to high uric acid, neurologic problems, and characteristic self-injurious behavior.

\medterm{Lesion} Any localized abnormal change in tissue caused by disease, injury, inflammation, infection, or a tumor. A lesion may be visible, felt, or found only on imaging or microscopic examination.
\medterm{Lesion Identification} The process of locating and describing an abnormal area in the body using examination, imaging, endoscopy, laboratory tests, or tissue sampling.
\medterm{Lesions On The Brain (Brain Lesions (Lesions On The Brain))} Alternate index label for Brain Lesions (Lesions on the Brain).

\medterm{Lesser Curvature} The shorter, inward curve along the upper and inner edge of the stomach. It is attached to a fold of tissue leading to the liver and contains blood vessels that supply the stomach.

\medterm{Lesser Trochanter} A small bony projection on the inner upper thighbone where the iliopsoas muscle attaches. It helps bend the hip; injury may occur from a sudden muscle pull in adolescents or, in adults, from weakened bone.


\medterm{Let-Down Reflex} Release of breast milk when oxytocin makes the small muscles around milk-producing glands contract. Suckling, breast stimulation, or seeing or hearing a baby can trigger it, although stress may delay it.
\medterm{Lethal} Able to cause death. In medicine, the word may describe a disease, injury, dose of a substance, or genetic change that prevents survival.

\medterm{Lethal-Means Counseling} A preventive intervention that reduces access to firearms, medications, toxic
substances, or other methods of self-harm during periods of elevated risk. It is
conducted collaboratively and is an important component of suicide prevention.

\medterm{Lethargy} Lethargy is reduced energy, alertness, or responsiveness; it may result from sleep loss, infection, metabolic disease, medicines, depression, or neurologic illness, and sudden marked lethargy requires urgent assessment.

\medterm{Letrozole} Letrozole is an aromatase inhibitor that lowers oestrogen production and is used mainly for hormone-receptor-positive breast cancer and ovulation induction; hot flushes, arthralgia, bone loss, and fetal harm are important considerations.



\medterm{Leucocyte} Another spelling of leukocyte, meaning a white blood cell. Leukocytes help defend against infection and other threats; major groups include neutrophils, lymphocytes, monocytes, eosinophils, and basophils.
\medterm{Leucocyte (White Blood Cell)} A blood cell that helps defend the body against infections, foreign substances, and abnormal cells; major types include lymphocytes, neutrophils, monocytes, eosinophils, and basophils.

\synonyms
White Blood Cell

\medterm{Leucocytosis} A higher-than-usual number of white blood cells in the blood. It may occur with infection, inflammation, physical stress, medicines, smoking, or blood and bone-marrow disorders, so the type of cell and the person's symptoms help find the cause.
\medterm{Leucoderma} Loss or marked reduction of skin pigmentation, often used as a descriptive term for depigmented patches.
\medterm{Leucoderma (Vitiligo)} Loss of skin pigment producing sharply defined pale or white patches, usually because melanocytes have been destroyed or have stopped functioning.

\synonyms
Vitiligo

\medterm{Leucomalacia} Softening or loss of tissue after injury, poor blood flow, infection, or inflammation. In newborns, the term often refers to damage of the brain’s white matter near the fluid spaces, which may affect movement or development.
\medterm{Leucopenia} A lower-than-usual number of white blood cells. It may result from medicines, infections, immune disorders, nutritional deficiency, or bone-marrow disease and can increase the risk of infection, especially when neutrophils are low.

\medterm{Leucoplakia} A persistent white patch on a moist body lining, often inside the mouth, that cannot be wiped away. It may result from irritation and sometimes contains precancerous changes, so it should be examined.
\medterm{Leucopoiesis} The production and maturation of white blood cells, mainly in the bone marrow. It is regulated by growth factors and increases or changes during infection, inflammation, and some blood disorders.
\medterm{Leucorrhoea (Leukorrhoea)} A whitish vaginal discharge that may be physiological or may result from infection, inflammation, or another genital disorder.

\synonyms
Leukorrhoea

\medterm{Leucovorin} Leucovorin is a folate medicine used to reduce the toxicity of methotrexate, enhance some chemotherapy regimens, and treat certain folate-related problems.

\medterm{Leucovorin Calcium} Leucovorin calcium is the calcium salt of leucovorin, used as rescue therapy after methotrexate and in selected colorectal-cancer treatments.

\medterm{Leukaemia} A cancer of blood-forming tissues in which abnormal white blood cells multiply uncontrollably and crowd out normal cells needed for immunity, clotting, and oxygen transport.
\medterm{Leukapheresis} A blood-filtering procedure that removes excess white blood cells and returns the remaining blood. It may rapidly reduce very high counts causing symptoms in leukemia, but it is temporary and does not replace treatment of the underlying disease.

\medterm{Leukemia} Leukemia is cancer of blood-forming cells that disrupts normal marrow production and may cause anemia, infection, bleeding, abnormal blood counts, and organ infiltration.
Leukemias are classified by the cell lineage and by clinical course, including acute and
chronic forms.

\synonyms
Blood Cancer (Leukemia)
Leukemia, General


\medterm{Leukemia Cutis} A condition in which leukemia cells collect in the skin, causing firm red, purple, or brown bumps, nodules, or patches. It often accompanies acute myeloid leukemia and requires assessment of the underlying blood cancer.

\medterm{Leukemia Genetics} Changes in genes or chromosomes in leukemia cells help identify the leukemia type, estimate its likely course, and choose treatment. The meaning of each change depends on the person’s specific leukemia.

\medterm{Leukemia, Acute Lymphocytic} Alternate index label; see \textit{Acute lymphocytic leukemia}.

\medterm{Leukemia, Acute Myelogenous} Alternate index label; see \textit{Acute myelogenous leukemia}.

\medterm{Leukemia, Chronic Lymphocytic} Alternate index label; see \textit{Chronic lymphocytic leukemia}.

\medterm{Leukemia, Chronic Myelogenous} Alternate index label; see \textit{Chronic myelogenous leukemia}.

\medterm{Leukemia, General} Alternate index label; see \textit{Leukemia}.

\medterm{Leukemia, Hairy Cell} Alternate index label; see \textit{Hairy cell leukemia}.

\medterm{Leukemic Infiltration Of The Gingiva} Swelling of the gums caused by leukemia cells collecting in gum tissue. The gums may become red, soft, enlarged, and prone to bleeding, especially in some acute myeloid leukemias.

\medterm{Leukemoid Reaction} A marked reactive increase in the white-cell count, usually with prominent neutrophilia
and immature granulocytes, caused by severe infection, inflammation, malignancy, or
tissue injury. It resembles leukemia but is not a clonal hematologic malignancy.

\medterm{Leukoaraiosis} White-matter changes seen on brain scans, often related to aging and long-term small-vessel disease such as high blood pressure. Extensive changes may be associated with slower thinking, walking problems, or stroke risk.

\medterm{Leukocyte} A white blood cell that helps defend the body against infection and other threats. Main types include neutrophils, lymphocytes, monocytes, eosinophils, and basophils; unusually high or low numbers can reflect infection, inflammation, medicines, or bone-marrow disease.

\medterm{Leukocyte Adhesion Deficiency} A rare inherited immune disorder in which white blood cells cannot attach to blood-vessel walls and reach infected tissue normally. It can cause delayed umbilical-cord separation, repeated skin infections with little pus, poor wound healing, gum disease, and dangerously severe infection in infants.

\medterm{Leukocyte Alkaline Phosphatase} A measure of an enzyme in certain white blood cells. It was formerly used to help distinguish a reactive response to infection from chronic myeloid leukemia, but newer genetic and blood tests are usually more informative.


\medterm{Leukocyte Disorder} A disorder involving the number, function, or development of white blood cells, including leukopenia, leukocytosis, immune deficiency, and malignant leukocyte proliferation.


\medterm{Leukocyte Esterase Urine Test} A urine dipstick test detecting leukocyte esterase released by white blood cells; a positive result supports pyuria but does not by itself diagnose urinary tract infection.


\medterm{Leukocyte Transfusion} Transfusion of donated white blood cells in rare, carefully selected situations involving life-threatening infection and severe inability to produce neutrophils. Benefits are uncertain and risks include fever, lung injury, immune reactions, and infection transmission.

\medterm{Leukocytes} White blood cells that defend the body against infections, abnormal cells, and other threats. Major types include neutrophils, lymphocytes, monocytes, eosinophils, and basophils, each with different immune functions.

\medterm{Leukocytoclastic Vasculitis} Inflammation of small blood vessels, usually in the skin, often causing raised purple or red spots that can be felt. Medicines, infections, or immune diseases may trigger it, and internal organs can sometimes be affected.

\medterm{Leukocytosis} A higher-than-usual number of white blood cells. It may occur with infection, inflammation, stress, smoking, medicines, or blood and bone-marrow disorders; the specific cell type and symptoms help identify the cause.

\medterm{Leukodystrophy} A group of inherited disorders that damage myelin, the protective covering around nerves in the brain and spinal cord. They can cause progressive problems with movement, thinking, vision, hearing, or other nervous-system functions.

\medterm{Leukoencephalopathy} Leukoencephalopathy is disease of the brain’s white matter caused by inherited, infectious, toxic, vascular, metabolic, inflammatory, or other processes.

\medterm{Leukonychia} White areas or bands in one or more nails. Causes include minor injury to the nail, inherited traits, infection, medicines, or illness; the pattern and whether it grows out with the nail help determine the cause.

\medterm{Leukonychia Striata} Longitudinal white bands running along the nail plate, usually caused by changes in the nail matrix. They may be inherited or follow repeated minor trauma, skin disease, or some medicines; persistent or widespread changes should be medically assessed.

\medterm{Leukonychia Totalis} Uniform white discoloration involving most or all of one or more nail plates. It is often inherited but can occasionally be associated with illness or medication; persistent changes affecting many nails should be evaluated with the medical history and other findings.

\medterm{Leukopenia} An abnormally low white-blood-cell count, which may increase susceptibility to infection depending on the cells affected and the severity of the reduction.

\medterm{Leukoplakia} A persistent white patch on the lining of the mouth or another mucosal surface that cannot be scraped away and cannot be explained by another condition. Some patches can develop into cancer, so persistent lesions need examination and often a biopsy.

\synonyms
Smokers Keratosis (Leukoplakia)

\medterm{Leukopoiesis} The process by which the bone marrow makes and develops white blood cells. Production increases or changes during infection, inflammation, immune disorders, and some blood diseases.

\medterm{Leukorrhea} Whitish vaginal discharge that may be normal, especially during hormonal changes, or may result from infection or inflammation. Foul odor, itching, pain, bleeding, or a new change requires evaluation.
\medterm{Leukostasis} A medical emergency in which very high numbers of abnormal leukocytes impair microvascular blood flow, potentially causing respiratory distress, neurologic symptoms, organ ischemia, or hemorrhage.

\medterm{Leukotriene} A leukotriene is an inflammatory lipid mediator made from arachidonic acid that contributes to bronchoconstriction, mucus, vascular leak, and immune-cell recruitment.

\medterm{Leukotriene Blockers} Medicines that reduce the effects or production of leukotrienes, substances involved in inflammation and airway narrowing. They help prevent some asthma or allergy symptoms but do not usually relieve a sudden asthma attack.

\synonyms
Leukotriene antagonists


\medterm{Leukotriene Modifiers} Medicines that reduce the production or effects of leukotrienes, inflammatory chemicals narrowing airways and causing allergy symptoms. They help prevent selected asthma or allergy problems but do not replace emergency rescue treatment.

\medterm{Leukotriene Receptor} A cell-surface protein activated by leukotrienes, inflammatory chemicals that can narrow airways, increase mucus, and promote swelling. Blocking these receptors can help prevent selected asthma and allergy symptoms.

\medterm{Leukotrienes} Chemical signals released during inflammation that can narrow airways, increase mucus, and cause swelling. They contribute to asthma, allergies, and some other inflammatory conditions.

\medterm{Leuprolide} Leuprolide is a gonadotropin-releasing hormone agonist used to suppress sex-hormone production in conditions such as prostate cancer, endometriosis, and fibroids.

\medterm{Leuprolide Acetate} Leuprolide acetate is the salt form of leuprolide, used to suppress ovarian or testicular hormone production for selected medical conditions.

\medterm{Levalbuterol} An inhaled short-acting bronchodilator that stimulates beta-2 receptors and relaxes airway muscle. It relieves wheezing and shortness of breath but may cause tremor, nervousness, headache, or a fast heartbeat.
\medterm{Levamisole} Levamisole is an antiparasitic and immune-modulating drug used selectively; serious low white-cell counts and vasculitic reactions can occur.

\medterm{Levator Ani} A group of pelvic-floor muscles that supports the bladder, bowel, and reproductive organs and helps control urination and bowel movements. Weakness or injury may contribute to pelvic organ prolapse or incontinence.

\medterm{Levetiracetam} Levetiracetam is a broad-spectrum antiseizure medicine that binds synaptic vesicle protein SV2A; it treats focal and generalized seizures and may cause somnolence, dizziness, or behavioral and mood changes.

\medterm{Levobunolol} Levobunolol is a beta-blocker eye drop used to lower pressure inside the eye in glaucoma or ocular hypertension.

\medterm{Levocardia} Levocardia is a normal left-sided position of the heart; it may occur with normal organ arrangement or with other congenital laterality abnormalities.

\medterm{Levocarnitine} Levocarnitine is the biologically active form of carnitine used to treat selected carnitine deficiencies and help transport fatty acids into mitochondria.

\medterm{Levodopa} Levodopa is a dopamine precursor that crosses the blood--brain barrier and is converted to dopamine, usually with carbidopa or benserazide, to treat Parkinson disease; dyskinesias, hallucinations, nausea, and motor fluctuations can develop.

\medterm{Levofloxacin} Levofloxacin is a fluoroquinolone antibiotic that inhibits bacterial DNA gyrase and topoisomerase IV; tendon injury, peripheral neuropathy, QT prolongation, dysglycaemia, and central-nervous-system effects are important class risks.

\medterm{Levoleucovorin} Levoleucovorin is the active isomer of leucovorin, used as methotrexate rescue and with selected chemotherapy regimens.

\medterm{Levoleucovorin Calcium} Levoleucovorin calcium is the calcium salt of levoleucovorin, used to reduce methotrexate toxicity and support selected cancer treatments.

\medterm{Levomepromazine} A phenothiazine medicine with antipsychotic, sedative, antiemetic, and analgesic properties, used in selected palliative-care and psychiatric settings.
\medterm{Levonorgestrel} A progestin hormone used in birth-control pills, hormonal intrauterine devices, and emergency contraception to prevent or delay ovulation and alter cervical mucus.

\medterm{Levorphanol} Levorphanol is a potent opioid analgesic that can relieve severe pain but also cause sedation, dependence, and respiratory depression.

\medterm{Levothyroxine} A synthetic form of thyroxine (T4) used as thyroid hormone replacement in hypothyroidism and in selected other thyroid conditions; dose is individualized and monitored with laboratory tests.

\medterm{Levothyroxine Sodium} A manufactured form of thyroid hormone used to replace hormone missing in hypothyroidism. The dose is adjusted using blood tests, and it is usually taken consistently on an empty stomach.

\synonyms
Levothyroxine

\medterm{Levulose} An older name for fructose, a simple sugar found in fruit, honey, and many sweeteners.
\medterm{Lewy Body} A Lewy body is an abnormal neuronal inclusion containing alpha-synuclein, characteristic of Parkinson disease and Lewy-body dementia.

\medterm{Lewy Body Dementia} A progressive brain disorder caused by abnormal alpha-synuclein protein deposits called Lewy bodies. It causes dementia with changing attention or alertness, recurrent visual hallucinations, Parkinson-like movement problems, sleep disorders such as acting out dreams, and sometimes automatic body-function problems such as constipation or dizziness on standing. Some people are very sensitive to antipsychotic medicines, so treatment must be individualized.

\synonyms
Dementia, Lewy Body

\medterm{Lewy Body Disease} A progressive brain disorder caused by abnormal protein deposits in nerve cells. It may cause changing attention, visual hallucinations, movement problems, stiffness, sleep-related behaviors, and gradual difficulty with thinking and daily activities.


\medterm{Leydig Cell} A cell between the testicular tubules that produces testosterone, mainly in response to luteinizing hormone.

\medterm{Leydig Cell Testicular Tumor} A usually rare sex-cord stromal tumor arising from testosterone-producing Leydig cells in the testis; it may cause a testicular mass or hormone-related signs.

\medterm{Leydig Cell Tumor} A testicular sex cord-stromal tumor that secretes testosterone or estrogen, causing
precocious puberty or gynecomastia; most are benign and treated by orchidectomy.

\medterm{Leydig Cells} Cells found between the sperm-making tubes in the testes. In response to luteinising hormone, they make testosterone, which supports male sexual development, sperm production, sex drive, and muscle and bone health.
\medterm{Löfgren Syndrome} A sudden form of sarcoidosis causing swollen lymph nodes near the lungs, tender red lumps on the shins, and painful swollen ankles. It often improves over time, but medical follow-up is needed to monitor symptoms and organ involvement.

\medterm{LGE} Alternate name; see \textit{Late Gadolinium Enhancement}.

\medterm{LH} Luteinising hormone, a hormone released by the pituitary gland. In people with ovaries it helps trigger ovulation and supports hormone production; in people with testes it stimulates testosterone production. Blood tests can help assess reproductive or pituitary problems.

\synonyms
Luteinizing hormone

\medterm{LH Response To GnRH Blood Test} A hormone test that measures luteinizing hormone before and after a medicine that stimulates its release. It helps assess delayed or early puberty and problems involving the brain, pituitary gland, ovaries, or testes.

\medterm{Libman-Sacks Endocarditis} Noninfectious inflammation of heart valves that can occur with lupus or antiphospholipid syndrome. Small deposits may form on the valves, sometimes causing leakage or increasing the risk of blood clots and stroke.

\synonyms
Verrucous endocarditis; Nonbacterial verrucous endocarditis

\medterm{Li-Fraumeni Syndrome} An inherited condition that greatly increases the chance of developing several cancers, often at a young age. Cancers may include breast cancer, sarcoma, brain tumors, adrenal cancer, or leukemia; specialized screening is usually advised.

\medterm{LI-RADS} A standardized system for reporting liver-imaging findings in people at increased risk of liver cancer, especially those with cirrhosis. Its categories estimate how likely a finding is to be cancer and help guide follow-up or biopsy.


\medterm{Libido} A person's sexual desire or interest. It can change with age, relationships, mood, stress, sleep, illness, pregnancy, menopause, hormone levels, and medicines; a change is important mainly when it causes distress or affects a person's wellbeing.

\medterm{Libido And Sexual Performance} Sexual desire and sexual function can be influenced by hormones, blood flow, nerves, medicines, illness, mood, stress, sleep, and relationships. A persistent or distressing change deserves an assessment for treatable causes.

\medterm{Lice} Small insects that live on hair or clothing and feed on blood. Head, body, and pubic lice cause itching and may leave eggs attached to hair. They spread mainly through close contact and are not a sign of poor hygiene.

\medterm{Lice Infestation} Infestation of the hair or skin by small insects that feed on blood. Head, body, and pubic lice spread mainly through close contact and cause itching or visible eggs attached to hair. Treatment uses an approved lice-killing product and fine combing; close contacts and clothing may need attention.

\medterm{Lice Infestations} Pediculosis is infestation by head, body, or pubic lice. Head lice cause scalp pruritus
and nits attached to hair shafts; body lice live mainly in clothing; pubic lice affect
coarse hair-bearing regions.


\medterm{Lice, Body} Alternate index label; see \textit{Body lice}.

\medterm{Licensed Clinical Social Worker} A social worker licensed to provide clinical assessment and psychotherapy or other mental-health services within the scope of local law and training.

\medterm{Lichen} A descriptive term for a firm, flat-topped, raised skin lesion or for a skin pattern resembling lichen planus. It describes appearance rather than a single disease or cause.

\medterm{Lichen Planopilaris} A lymphocytic inflammatory scalp disorder that attacks hair follicles and can cause scarring alopecia, perifollicular scale, redness, and permanent hair loss.

\medterm{Lichen Planus} An inflammatory condition that can affect the skin, mouth, genitals, nails, or hair. It often causes itchy purple or reddish bumps and may produce sore or white-patterned patches in the mouth; the cause is often not known.

\synonyms
Lichen Ruber Planus (Lichen Planus)

\medterm{Lichen Planus, Oral} Alternate index label; see \textit{Oral lichen planus}.

\medterm{Lichen Ruber Planus (Lichen Planus)} Alternate index label for Lichen Planus.

\medterm{Lichen Sclerosus} A long-term inflammatory skin condition that usually affects the genital or anal area, causing pale, thin, itchy, easily injured skin. Treatment relieves symptoms and lowers scarring; persistent sores need review because cancer risk is slightly increased.

\synonyms
White Spot Disease

\medterm{Lichen Simplex Chronicus} A thickened, leathery plaque caused by chronic scratching and rubbing (itch-scratch
cycle). It commonly affects the neck, ankles, lower legs, and genital area. It is a form
of neurodermatitis associated with stress and anxiety.

\medterm{Lichenification} Thick, leathery skin with more visible skin lines caused by repeated scratching or rubbing. It often develops in long-lasting eczema, allergies, or other itchy skin conditions.

\medterm{Lichenified} Lichenified skin is thickened and leathery from chronic rubbing or scratching, commonly occurring in longstanding eczema or itch disorders.

\medterm{Lichenoid Drug Eruption} A lichenoid inflammatory skin reaction caused by a medicine, often resembling lichen planus but with a broader distribution or different clinical and histologic features; improvement may follow withdrawal of the culprit drug.

\medterm{Liddle Syndrome} A rare inherited cause of early high blood pressure in which the kidneys retain too much sodium and lose too much potassium. It may begin in childhood or young adulthood and usually requires specialist treatment.

\medterm{Liddle Syndrome Treatment} Treatment of Liddle syndrome lowers blood pressure and the body’s excessive sodium retention, usually with a medicine that blocks the overactive kidney channel. Potassium and kidney function are monitored, and relatives may need testing.

\medterm{Lidocaine} Lidocaine is a local anaesthetic that blocks voltage-gated sodium channels and, intravenously, can treat selected ventricular arrhythmias; excessive exposure may cause seizures, hypotension, arrhythmia, or cardiac arrest.

\medterm{LieberküHn'S Crypts} Tubular glands in the intestinal lining between villi that contain stem, absorptive, and secretory cells. They renew the lining, produce intestinal substances, and help maintain the gut’s protective barrier.
\medterm{Life Expectancy} The average number of additional years a person of a specified age is expected to live under current mortality conditions.


\medterm{Lifetime Risk} The probability that a person will develop a specified disease or experience an outcome during the remainder of their life or over a defined lifetime period.

\medterm{Ligament} A dense, fibrous connective tissue band that connects bone to bone, providing passive
joint stability and guiding joint motion. Ligaments are composed predominantly of type I
collagen fibers (parallel orientation, crimped pattern) and elastin fibers (10-20\% in
some ligaments), arranged in bundles.

\medterm{Ligand} Any molecule that binds specifically to a biological receptor, macromolecule, or active
site. In pharmacology and physiology, ligands include hormones, neurotransmitters,
drugs, and toxins that bind receptors to initiate (agonist), block (antagonist), or
modulate signaling.

\medterm{Ligase} An enzyme that joins two molecules, often using energy from ATP. Ligases help repair DNA, copy genetic material, and build or maintain other important molecules.

\medterm{Ligation} The surgical process of tying off a blood vessel or duct using a suture or clip to
prevent bleeding or block fluid flow.

\medterm{Ligation (Vascular)} Surgical tying, clipping, or sealing of a blood vessel to stop bleeding or redirect blood flow. It may be used during surgery or to treat selected abnormal vessels.
\synonyms
Vascular ligation

\medterm{Ligature} Any cord, wire, rope, fabric, or flexible object used to constrict the neck, bind limbs,
suspend the body, or inflict injury. The resulting pale groove-shaped abrasion may be
horizontal with petechiae in homicidal strangulation, or oblique and upward in hanging.

\medterm{Light} Visible electromagnetic radiation; in medicine it is used for examination, phototherapy, imaging, and selected treatments.
\medterm{Light Coagulation} Light coagulation uses intense focused light, usually a laser, to heat and seal tissue or blood vessels.

\medterm{Light Cycle} A light cycle is the recurring pattern of light and darkness that helps regulate circadian rhythms and biological activity.

\medterm{Light Microscope} An instrument that uses visible light and lenses to magnify cells, tissues, microorganisms, or other small structures. It is widely used in clinical laboratories, pathology, microbiology, and research.

\medterm{Light Microscopy} Examination of cells, tissues, microorganisms, or other specimens using visible light and a microscope. Different stains and preparation methods reveal structure, infection, inflammation, blood cells, or abnormal growth.

\medterm{Light Reflex} Automatic narrowing of the pupil when light enters the eye. It tests the retina, optic nerve, brain pathways, and nerve that controls the pupil; an unequal or absent response may indicate disease.
\medterm{Light Scattering} Deflection of light when it meets particles, cells, molecules, or irregular structures. Measuring the scattered light can help estimate particle size or concentration in laboratory tests and can affect the appearance of tissues on imaging.

\medterm{Light Therapy} Treatment that uses selected wavelengths or intensity of light for conditions such as some skin disorders, seasonal depression, or neonatal jaundice.

\medterm{Light's Criteria} A set of laboratory comparisons used to decide whether fluid around the lung is caused by inflammation or disease in the chest rather than by pressure or fluid imbalance elsewhere. Protein and enzyme levels in the fluid are compared with blood levels; the result helps guide further testing and treatment.

\medterm{Light-Near Dissociation} A pupil abnormality in which the pupils constrict when focusing on a nearby object but respond poorly or not at all to light. It can occur with certain nerve disorders, infections, diabetes, or midbrain disease.

\medterm{Lightening} In late pregnancy, descent of the fetus into the pelvis before labor, which may reduce upper-abdominal pressure but increase pelvic pressure and urinary frequency.

\medterm{Lighter Fluid Poisoning} Harm from swallowing or breathing lighter fluid, a hydrocarbon. The greatest danger is liquid entering the lungs, causing chemical pneumonia, cough, or breathing difficulty; do not induce vomiting and contact poison control immediately.

\medterm{Lightheadedness} A feeling of near-fainting, floating, or unsteadiness without necessarily perceiving that the surroundings are spinning; causes include dehydration, low blood pressure, anemia, medications, and heart or neurologic conditions.

\medterm{Lightning Pains} Brief, stabbing, electric-shock-like pains caused by irritated or damaged nerves. They classically occur in tabes dorsalis but may also accompany shingles, nerve compression, diabetes, or other neurological disorders.
\medterm{Lignocaine} The international nonproprietary name for lidocaine, a local anesthetic and antiarrhythmic drug.
\medterm{Likelihood Ratio} A measure of how strongly a test result supports or argues against a disease. A value above one supports disease, a value below one argues against it, and larger changes are more informative.

\medterm{Lilliputian Hallucination} A visual hallucination of unusually small people, animals, or objects; it can occur with neurologic, psychiatric, toxic, infectious, or sensory conditions.

\medterm{Lily Of The Valley Poisoning} Poisoning from eating lily-of-the-valley plant parts, which contain substances that can disrupt the heartbeat. Nausea, vomiting, dizziness, slow or irregular pulse, confusion, or fainting require immediate emergency care.

\medterm{Limb} An arm or leg, including its bones, joints, muscles, tendons, nerves, blood vessels, and skin. Limb problems may affect movement, sensation, strength, circulation, or function.

\medterm{Limb Bud} A limb bud is an embryonic outgrowth that develops into an arm or leg through coordinated signaling, growth, and tissue differentiation.

\medterm{Limb Development} Limb development is embryonic formation and growth of the bones, joints, muscles, nerves, vessels, and skin of the arms and legs.

\medterm{Limb Reduction Defects} Birth defects in which part or all of an arm or
leg does not form normally. The missing or shortened part may involve a bone,
joint, hand, foot, or the whole limb.

\medterm{Limb Fracture} A break in a bone of an upper or lower limb, assessed according to location, displacement, open or closed status, joint involvement, and associated neurovascular injury.

\medterm{Limb Girdle} A limb girdle is the skeletal connection of a limb to the trunk, such as the shoulder or pelvic girdle.

\medterm{Limb Injury} Trauma affecting an arm or leg, including injury to bone, joint, muscle, tendon, ligament, nerve, blood vessel, or skin.

\medterm{Limb Pain} Pain arising from an arm or leg, with causes including trauma, nerve disease, vascular insufficiency, infection, inflammation, and musculoskeletal disorders.

\medterm{Limb Plethysmography} A noninvasive test measuring changes in limb volume or blood flow, used to evaluate arterial or venous obstruction and sometimes to assess vascular responses.

\medterm{Limb Prosthesis} A limb prosthesis is an artificial device replacing part or all of an arm or leg to restore mobility, function, appearance, or independence.

\medterm{Limb-Girdle Muscular Dystrophies} Limb-girdle muscular dystrophies are inherited muscle diseases causing progressive weakness predominantly around the hips and shoulders, with variable cardiac and respiratory involvement.

\medterm{Limbic Encephalitis} Inflammation of the limbic system, often presenting subacutely with memory impairment,
seizures, psychiatric or behavioral change, and sleep disturbance. It may be autoimmune,
paraneoplastic, infectious, or of unknown cause.

\medterm{Limbic Lobe} A curved region on the inner surface of the brain involved in memory, emotion, motivation, and some responses controlled by the nervous system.
\medterm{Limbic System} A group of connected brain structures involved in memory, emotion, motivation, learning, and some automatic body responses, such as changes in heart rate during stress.

\medterm{Limbs} The arms and legs, including their bones, joints, muscles, nerves, blood vessels, and skin. Limb problems may affect strength, sensation, movement, circulation, balance, daily activities, or independence.
\medterm{Limbus} The border where the clear cornea meets the white outer coat of the eye. It contains cells that help renew the corneal surface and may be affected by injury, inflammation, or disease.

\synonyms
Corneoscleral limbus

\medterm{Limbus Corneae} The corneal limbus is the circular border where the transparent cornea meets the white sclera and conjunctiva, containing stem cells important for corneal repair.

\medterm{Limited Range Of Motion} Limited range of motion is reduced ability to move a joint through its normal extent because of pain, stiffness, weakness, swelling, or structural change.

\medterm{Linagliptin} Linagliptin is a DPP-4 inhibitor used with diet and exercise to improve blood glucose control in type 2 diabetes.

\medterm{Lincomycin} Lincomycin is a lincosamide antibiotic used for selected serious bacterial infections; it can cause antibiotic-associated diarrhea and colitis.

\medterm{Lindane} Lindane is an organochlorine insecticide formerly used for lice or scabies; neurotoxicity and environmental persistence limit its use.

\medterm{Line, Central} A catheter whose tip lies in a large central vein, used to deliver medicines or fluids, monitor central venous pressure, draw blood, or provide specialized nutrition.

\medterm{Line, Central Venous} A central venous catheter inserted into a large vein such as the internal jugular, subclavian, or femoral vein and advanced toward the superior or inferior vena cava.

\medterm{Linea} A line or linear anatomical marking, often used in names such as linea alba or linea nigra.
\medterm{Linea Aspera} Prominent longitudinal posterior femoral shaft ridge from converging muscle attachment
lines. Insertion for adductor magnus, vastus medialis/lateralis, and short head of
biceps femoris. Thickest femoral cortex, resisting bending forces. Diverges distally
into supracondylar lines.

\medterm{Linea Nigra} A dark brown or black vertical line appearing on the abdomen during pregnancy, extending
from the pubic symphysis to the umbilicus or beyond. Caused by hormonal stimulation of
melanocytes, resulting in hyperpigmentation. More prominent in darker-skinned women.

\medterm{Linear} Arranged in a straight or nearly straight line, describing the shape or distribution of a skin lesion or eruption.

\medterm{Linear Accelerator} A machine accelerating electrons to produce high-energy X-rays or electron beams that precisely deliver radiation treatment to selected tumors.
\medterm{Linear Energy Transfer} The amount of energy that radiation deposits as it travels through tissue. Radiation with higher transfer deposits energy more densely and may cause more biological damage for the same total dose.

\medterm{Linear Fracture} A narrow break that runs along the length of a bone without moving the pieces far apart or breaking the bone into many fragments. Treatment depends on the bone, the injury's stability, and whether nearby nerves, blood vessels, or organs are affected.
\medterm{Linear IgA Bullous Dermatosis} An autoimmune blistering disease caused by linear deposition of IgA at the
basement membrane zone. It produces tense blisters, sometimes in annular or "string of pearls" patterns, and may be
idiopathic or drug-induced.

\medterm{Lined Canister} A container with a protective inner lining used to collect or store a medical substance, specimen, waste, or device component. The lining helps prevent leakage, contamination, chemical reaction, or damage to the container; its exact use depends on the equipment.

\medterm{Linezolid} Linezolid is an oxazolidinone antibiotic that blocks formation of the bacterial 70S initiation complex and treats resistant Gram-positive infections; thrombocytopenia, neuropathy, lactic acidosis, and serotonin interactions require caution.

\medterm{Lingua} The anatomical Latin term for the tongue, a muscular organ used for chewing, swallowing, speaking, and tasting. The front two-thirds lies in the mouth and the back third lies in the throat; related terms include sublingual, meaning under the tongue.
\medterm{Lingua Villosa Nigra} A benign condition characterized by elongation and dark pigmentation of the filiform
papillae of the tongue, often associated with tobacco use, poor oral hygiene, or
antimicrobial therapy.

\medterm{Lingual} Relating to the tongue or the tongue side of a body structure. Lingual muscles, nerves, blood vessels, and surfaces help with taste, speech, swallowing, and movement of food.

\medterm{Lingual Artery} A branch of the external carotid artery supplying the tongue, floor of mouth, and
sublingual structures. It arises anterior to the hypoglossal nerve and courses deep to
the hyoglossus muscle. Lingual artery aneurysm or thrombosis is rare but can cause
life-threatening hemorrhage or tongue ischemia.

\medterm{Lingual Thyroid} A lingual thyroid is thyroid tissue at the base of the tongue caused by failure of normal thyroid descent and may obstruct swallowing or breathing or cause hypothyroidism.

\medterm{Lingual Tonsil} The lingual tonsil is lymphoid tissue at the base of the tongue that contributes to immune surveillance of the throat.

\medterm{Lingual Tonsils} Collections of lymphoid tissue at the base of the tongue.

\medterm{Liniment} A liquid or semiliquid preparation rubbed onto the skin, usually to provide temporary relief of muscle or joint discomfort. Some contain irritating or warming substances and should not be used on broken skin, near the eyes, or under tight coverings unless directed.

\medterm{Linitis Plastica} A form of stomach cancer that spreads widely through the stomach wall, making it thick and stiff. It may cause early fullness, weight loss, nausea, vomiting, or vague upper-abdominal discomfort.


\medterm{Linoleic Acid} An essential omega-6 polyunsaturated fatty acid obtained from food. The body uses it to make cell membranes and signaling molecules.
\medterm{Linolenic Acid} A polyunsaturated fatty acid. Alpha-linolenic acid is an essential omega-3 nutrient from plant foods, while other forms occur in animal tissues and have different roles.


\medterm{Liothyronine Sodium} A manufactured form of triiodothyronine, an active thyroid hormone, used in selected thyroid disorders. It acts quickly and can cause fast heartbeat, tremor, anxiety, bone loss, or heart strain if excessive.

\medterm{Lip} The lips are mobile folds of skin and mucosa containing muscle, blood vessels, and sensory nerves; they form the oral seal and assist speech, facial expression, feeding, and swallowing.

\medterm{Lip Cancer} Cancer developing in the tissue of the upper or lower lip. It may appear as a persistent sore, lump, crust, bleeding area, numbness, or color change and needs medical assessment.

\medterm{Lip Disease} A disorder of the lips, such as infection, inflammation, ulceration, trauma, pigmentation change, benign tumor, or malignancy.

\medterm{Lip Infection} Infection of the lip caused by bacteria, viruses, fungi, or an infected wound. It may cause pain, redness, swelling, blisters, cracks, pus, or fever; rapidly spreading swelling or breathing difficulty requires urgent care.
\medterm{Lip Moisturizer Poisoning} Swallowing a small amount of ordinary lip moisturizer usually causes mild stomach upset or loose stools. Products containing camphor, salicylates, or other medicines can cause serious poisoning and require poison-control advice.

\medterm{Lip Neoplasm} A lip neoplasm is an abnormal growth in the lip that may be benign or malignant; persistent sores or lumps require assessment.

\medterm{Lip Swelling} Enlargement of one or both lips caused by local injury, infection, inflammation, dental disease, medication, or angioedema. Sudden swelling of the lips with tongue or throat swelling, breathing difficulty, wheezing, faintness, or widespread hives may indicate a severe allergic reaction and requires emergency care.

\synonyms
Lip edema; lip swelling

\medterm{Lipase} An enzyme that breaks down fats into smaller molecules. Blood lipase testing is commonly used when pancreatitis is suspected, although an abnormal result is not specific in every case.

\medterm{Lipase Test} A lipase test measures a pancreatic digestive enzyme in blood and supports evaluation of acute pancreatitis and selected pancreatic or gastrointestinal disorders.

\medterm{Lipedema} A long-term disorder causing an unusual, usually symmetrical buildup of painful fat in the legs and sometimes arms. The feet and hands are usually spared, while easy bruising, tenderness, and swelling may occur.

\medterm{Lipemia Retinalis} A creamy pink-white appearance of the blood vessels at the back of the eye caused by extremely high blood triglycerides. It signals severe fat buildup in the blood and a high risk of sudden pancreatitis.

\medterm{Lipemic Index} A laboratory estimate of how cloudy a blood sample is because of excess fats. Marked cloudiness can interfere with some test results and may prompt review of the sample or repeat testing.

\medterm{Lipid} A fat or fat-like substance, including triglycerides, cholesterol, phospholipids, and fat-soluble vitamins. Lipids store energy, form cell membranes, help make hormones, and help the body absorb vitamins A, D, E, and K.

\medterm{Lipid Bilayer} The double layer of fat-like molecules forming the basic structure of cell membranes. It separates the cell from its surroundings and controls movement of substances.

\medterm{Lipid Disorder} An abnormal level or handling of fats in the blood or tissues, such as high LDL cholesterol, high triglycerides, or unusually low protective HDL cholesterol.
\medterm{Lipid Emulsion Therapy} Emergency intravenous treatment for severe poisoning from a local anesthetic or selected other fat-soluble substances. The fat may draw some drug away from the heart and brain and provide energy to the heart; it is given under medical supervision.

\medterm{Lipid Metabolism} The processes by which the body digests, transports, stores, makes, and breaks down fats for energy, cell membranes, and hormone production.
\medterm{Lipid Panel} A blood test measuring total cholesterol, LDL cholesterol, HDL cholesterol, and triglycerides. The results help estimate cardiovascular risk and guide lifestyle or medicine decisions; fasting may or may not be needed depending on the purpose.

\medterm{Lipid Profile} A blood test panel measuring lipids and lipoproteins, commonly including total cholesterol, LDL cholesterol, HDL cholesterol, and triglycerides.


\medterm{Lipid Profile Test} A blood test usually measuring total cholesterol, LDL cholesterol, HDL cholesterol, and triglycerides to assess atherosclerotic cardiovascular risk and monitor treatment.

\medterm{Lipid Transport} Movement of fats through blood and tissues by lipoproteins, carrier proteins, or other mechanisms for energy use, storage, membrane formation, and hormone production.

\medterm{Lipidosis} A disorder involving abnormal accumulation or metabolism of lipids in tissues, often due to inherited enzyme defects.
\medterm{Lipids} A group of mostly water-insoluble substances that includes fats, oils, fatty acids, phospholipids, and cholesterol. Lipids store energy, form cell membranes, help make hormones, and help the body absorb vitamins A, D, E, and K.

\medterm{Lipoatrophy} Loss of the layer of fat beneath the skin, causing dents or hollow areas. It can follow repeated injections, inflammation, autoimmune disease, HIV treatment, or inherited disorders and may alter medicine absorption at injection sites.

\medterm{Lipoblastoma} A usually benign tumor of immature adipose tissue occurring mainly in infants and young children; it may be localized or diffuse and can recur if incompletely removed.

\medterm{Lipodystrophy} Abnormal loss or distribution of body fat. It may be inherited or caused by medicines, metabolic disease, or repeated injections and can lead to insulin resistance, diabetes, high triglycerides, or changes in appearance.

\medterm{Lipodystrophy (Acquired Generalized And Inherited Lipodystrophy)} Alternate index label for Acquired Generalized and Inherited Lipodystrophy.

\medterm{Lipodystrophy Syndromes} Rare disorders in which body fat is absent from some areas, present from birth or lost later, and may collect excessively elsewhere. They can cause severe insulin resistance, diabetes, fatty liver, and very high triglycerides.

\medterm{Lipodystrophy, Intestinal} Alternate index label; see \textit{Whipple's disease}.

\medterm{Lipofuscin} Lipofuscin is a yellow-brown, non-degradable pigment made of oxidized lipid and protein residues that accumulates in long-lived cells with ageing and in some degenerative diseases.


\medterm{Lipogranulomatosis} Inflammation in which fat accumulates and triggers organized collections of immune cells called granulomas. It may affect skin, soft tissue, or internal organs and can result from injury, leakage of oily material, or metabolic disease.

\medterm{Lipohypertrophy} A firm or lumpy buildup of fat under the skin, often at repeated insulin-injection sites. Injecting into these areas can change insulin absorption, so injection locations should be rotated regularly.

\medterm{Lipoid Nephrosis} An older name for a kidney disorder, most often in children, in which the kidney filters leak large amounts of protein into the urine. It can cause swelling around the eyes, legs, or abdomen.

\medterm{Lipoid Pneumonia} Lung inflammation caused by oil or fat entering the airways or building up in lung tissue. It may follow inhalation or swallowing of oil-based products and can cause cough, breathlessness, fever, or chest changes on imaging.

\medterm{Lipoidosis} Abnormal buildup of fats within cells or tissues, usually because a metabolic pathway cannot process or transport them normally. It may be inherited or acquired and can damage organs such as the liver, nerves, or kidneys.

\medterm{Lipolysis} Breakdown of stored body fat into fatty acids and glycerol so they can be used for energy. Fasting and stress hormones increase this process, while insulin generally reduces it.

\medterm{Lipoma} A common noncancerous lump made of fat cells, usually soft, movable, painless, and slow-growing. It often occurs under the skin of the trunk, neck, or limbs; a deep, painful, hard, or rapidly growing lump needs assessment.


\medterm{Lipomatosis} A condition in which mature fat tissue grows in multiple or broad areas rather than forming one separate lump. It may be symmetrical, involve the trunk or limbs, and be difficult to remove when it extends between nearby structures.

\medterm{Lipomatous Neoplasm} A growth made partly or mainly of fat cells. It may be noncancerous, like a lipoma, or cancerous, like a liposarcoma; a deep, firm, painful, or rapidly enlarging lump needs specialist assessment.


\medterm{Lipophilicity} Lipophilicity is the tendency of a substance to dissolve in fats or oils rather than water, influencing drug absorption and distribution.

\medterm{Lipopolysaccharide} Lipopolysaccharide is the outer-membrane endotoxin of Gram-negative bacteria, consisting of lipid A, core polysaccharide, and O antigen; circulating lipid A can trigger cytokine release, fever, shock, and disseminated inflammation.

\medterm{Lipoprotein} A particle made of fat and protein that carries cholesterol and other fats through the blood. Different types transport dietary fat, triglycerides, or cholesterol; high LDL levels can contribute to artery disease.

\medterm{Lipoprotein Analysis} A blood test measuring cholesterol, triglycerides, and particles such as LDL and HDL to estimate cardiovascular risk and guide treatment.

\synonyms
Lipid panel; lipid profile

\medterm{Lipoprotein Lipase} An enzyme on capillary walls that breaks down triglycerides in chylomicrons and very-low-density lipoproteins so tissues can use or store the fatty acids. Its activity is regulated by nutritional and hormonal signals.

\medterm{Lipoprotein Metabolism} The body's production, transport, modification, and removal of particles that carry cholesterol and other fats through the bloodstream.

\medterm{Lipoprotein(A)} A cholesterol-carrying particle similar to LDL whose level is largely inherited. A high level can increase the long-term risk of heart attack, stroke, and narrowing of the aortic valve.

\synonyms
Lp(a)

\medterm{Liposarcoma} A malignant tumor of fat cells, usually arising in the deep soft tissues of the thigh or
retroperitoneum; well-differentiated, myxoid, and pleomorphic subtypes are recognized.

\medterm{Liposarcoma Of The Breast} A rare malignant mesenchymal tumor with adipocytic differentiation arising in breast soft tissue. Diagnosis requires histopathologic and molecular evaluation, and treatment is generally wide excision with sarcoma-focused care.

\medterm{Liposomal Amphotericin B} A liposomal formulation of amphotericin B used in the treatment of severe systemic
fungal infections, with reduced nephrotoxicity compared to conventional amphotericin B.
The liposomal formulation encapsulates amphotericin B within lipid vesicles, which
alters the drug's biodistribution and reduces renal tubular damage.

\medterm{Liposome} A liposome is a tiny vesicle made of a lipid bilayer that can carry medicines or other substances into or around the body.

\medterm{Liposuction} A surgical procedure that uses a cannula and suction to remove localized deposits of subcutaneous fat for body contouring. It is not a treatment for obesity or a substitute for diet and exercise. Risks may include bleeding, infection, contour irregularity, fluid accumulation or shifts, sensory changes, and anesthesia-related complications.

\synonyms
Lipoplasty

\medterm{Lips} The paired movable folds surrounding the opening of the mouth, involved in speech, facial expression, eating, and oral sensation.

\medterm{LipschüTz Bodies} An older term for nonspecific acute genital ulcers, often occurring in adolescents or young women after an infection.
\medterm{Liquefaction} Conversion of a solid or partly solid material into liquid, such as tissue breakdown during liquefactive necrosis or melting of a substance.
\medterm{Liquefactive Necrosis} A pattern of tissue death in which enzymes digest dead cells until the area becomes soft or liquid. It commonly occurs in brain injury and some infections and may leave a fluid-filled cavity.

\medterm{Liquid Biopsy} A test of blood, urine, or another body fluid for cancer cells or DNA, RNA, and other material released by a tumor. It may help identify genetic changes, guide treatment, or monitor response and recurrence, but it does not always replace a tissue biopsy.

\medterm{Liquid Chromatography} A laboratory method that separates dissolved substances as a liquid carries them through a material that interacts differently with each substance. It is used to identify or measure medicines, hormones, metabolites, proteins, and toxins.

\medterm{Liquid Diet} A diet consisting partly or entirely of liquids; clear-liquid and full-liquid diets differ in allowed foods and may be prescribed temporarily for procedures, illness, or swallowing or digestive problems.

\medterm{Liquid Incense} A scented liquid product that may contain solvents, oils, or other chemicals. Swallowing or inhaling it can cause irritation, drowsiness, vomiting, seizures, or lung injury; the product container is needed for poison-control assessment.

\medterm{Liquid Nitrogen} Liquid nitrogen is extremely cold nitrogen used to freeze tissue in cryotherapy; contact can cause frostbite and its gas can displace oxygen.

\medterm{Liquor} A liquid; in clinical usage it may refer to cerebrospinal fluid or a medicinal preparation, depending on context.
\medterm{Liraglutide} Liraglutide is a GLP-1 receptor agonist used to improve glucose control in type 2 diabetes and, at some doses, support weight management.

\medterm{Lisch Nodule} A Lisch nodule is a harmless pigmented hamartoma of the iris, commonly associated with neurofibromatosis type 1.

\medterm{Lisch Nodules} Benign, dome-shaped melanocytic hamartomas on the iris, commonly seen in neurofibromatosis type 1.
They usually do not affect vision and support the diagnosis when interpreted with other
clinical findings.

\medterm{Lisinopril} Lisinopril is an angiotensin-converting-enzyme inhibitor used for hypertension, heart failure, and kidney protection in some diabetes; cough, hyperkalaemia, renal-function changes, angioedema, and fetal toxicity are important risks.

\medterm{Lissencephaly} A birth condition in which the brain’s surface has few or no normal folds because nerve cells did not move to their usual positions during development. It can cause seizures, severe developmental impairment, feeding problems, and movement difficulties.

\medterm{Listeria} Listeria is a genus of bacteria; L. monocytogenes can cause foodborne illness, bloodstream infection, and meningitis, especially in pregnancy, older adults, and people with weakened immunity.

\synonyms
Infection Listeria (Listeria)

\medterm{Listeria Infection} An infection caused by Listeria bacteria, usually acquired from contaminated food. It may cause fever, muscle aches, diarrhea, or vomiting; infection can be severe in pregnancy, newborns, older adults, and people with weakened immune systems.

\medterm{Listeria Monocytogenes} A gram-positive, facultatively anaerobic rod that causes listeriosis, a foodborne
infection. The organism is transmitted through contaminated food, particularly deli
meats, soft cheeses, and ready-to-eat foods. It can grow at refrigeration temperatures.

\medterm{Listeriosis} Infection with Listeria monocytogenes, usually acquired from contaminated food. It may cause fever and diarrhea or severe disease such as bloodstream infection or meningitis, especially in pregnancy, newborns, older adults, and people with weakened immunity.

\medterm{Listeriosis In Pregnancy} A foodborne infection that may cause only fever, tiredness, or muscle aches in the pregnant person but can spread to the fetus. It can cause miscarriage, stillbirth, premature birth, or severe infection in the newborn.


\medterm{Lithiasis} The formation or presence of stones (calculi) in an organ or duct. The term is often named for the site, such as nephrolithiasis in the kidney, ureterolithiasis in the ureter, or cholelithiasis in the gallbladder.

\medterm{Lithium} Lithium is a mood stabilizer used mainly for bipolar disorder and selected augmentation treatment; its narrow therapeutic range requires monitoring for tremor, thyroid or renal dysfunction, dehydration, drug interactions, and toxicity.

\medterm{Lithium Carbonate} A mood-stabilizing medicine used mainly to treat bipolar disorder and reduce the risk of manic or depressive episodes; serum levels, kidney function, thyroid function, and drug interactions require monitoring.

\medterm{Lithium Toxicity} Poisoning caused by too much lithium in the body, often after dehydration, illness, a dose change, or a medicine interaction. Nausea, diarrhea, worsening tremor, weakness, confusion, poor coordination, seizures, or reduced consciousness require urgent medical care.

\medterm{Lithotomy} A surgical procedure for removing a stone, or the position with the hips and knees flexed and legs supported.
\medterm{Lithotomy Position} A supine position in which the hips and knees are flexed and the legs are raised and supported to provide access to the pelvis and perineum.

\medterm{Lithotripsy} A treatment that breaks urinary or other stones into smaller pieces using shock waves, laser energy, or another device so the fragments can pass or be removed.

\medterm{Lithotripsy Probe} A device that fragments a stone using mechanical, ultrasonic, laser, or shock-wave energy during a
stone-removal procedure.

\medterm{Lithotripter} A device that breaks urinary, bile-duct, or other stones into smaller pieces using shock waves, laser energy, or another treatment method.
\medterm{Lithotripter Probe} A specialized probe delivering mechanical, electrical, or laser energy to break urinary or bile-duct stones into fragments that can be removed.

\medterm{Litmus} A dye that changes color with acidity or alkalinity and is used as a simple pH indicator.
\medterm{Little Area} The highly vascular region of the anterior nasal septum where several arteries meet and where most anterior
nosebleeds originate.

\medterm{Little'S Disease} An older name for spastic cerebral palsy, a condition caused by early brain injury or abnormal development that produces stiff muscles and movement difficulties, mainly in the legs. It is not progressive, although needs change over time.

\medterm{Live Attenuated Vaccines} Vaccines containing weakened forms of viruses or bacteria. They usually produce a strong, long-lasting immune response without causing illness in healthy people, but some are unsuitable for people with severely weakened immune systems or during certain pregnancies.

\synonyms
Live vaccines

\medterm{Live Birth} A live birth is delivery of an infant showing signs of life such as breathing, heartbeat, or movement, regardless of gestational age.

\medterm{Livedo Racemosa} A persistent, irregular, branching net-like pattern of reddish-blue skin caused by disease of small or medium blood vessels. It may occur with autoimmune, clotting, inflammatory, or other systemic disorders and needs evaluation.

\medterm{Livedo Reticularis} A net-like, violaceous discoloration of the skin caused by altered blood flow in small cutaneous vessels.
It may be temporary and cold-induced or associated with vascular, autoimmune, infectious,
or clotting disorders.

\medterm{Livedoid Vasculopathy} A chronic occlusive disorder of small dermal vessels causing painful purpuric lesions, ulcers, and porcelain-white stellate scars, usually on the lower legs and ankles.

\medterm{Liver} The largest internal organ, beneath the right ribs. It processes nutrients, makes bile to help digest fats, stores energy, helps regulate blood chemicals, and breaks down medicines and harmful substances.

\medterm{Liver (Anatomy And Function)} Alternate index label for Anatomy and Function.

\medterm{Liver Abscess} A pocket of pus inside the liver caused by bacteria, parasites, or another infection. It may cause fever, chills, right-sided upper-abdominal pain, or an enlarged liver and usually needs antibiotics with drainage or other treatment.

\medterm{Liver Biopsy} Percutaneous, transjugular, or laparoscopic sampling of liver tissue for histologic evaluation when noninvasive tests do not answer the clinical question. It may stage fibrosis or characterize inflammation, fatty liver disease, infiltrative disease, or a mass; complications include pain and bleeding.

\medterm{Liver Cancer} A malignant tumor arising in the liver. Hepatocellular carcinoma is the most common primary type in adults; risk factors include chronic liver disease, viral hepatitis, alcohol-related liver injury, and some toxins.

\synonyms
Cancer, Liver

\medterm{Liver Cancer - Hepatocellular Carcinoma} A cancer beginning in hepatocytes, the main liver cells, often associated with cirrhosis or chronic hepatitis.

\medterm{Liver Circulation} The liver receives oxygen-rich blood from the hepatic artery and nutrient-rich blood from the portal vein. Blood passes through liver tissue and leaves through the hepatic veins to the heart.

\medterm{Liver Cirrhosis} Permanent scarring of the liver after long-term injury, causing normal liver tissue to be replaced by stiff scar tissue. Causes include viral hepatitis, alcohol-related disease, fatty liver disease, autoimmune disease, and inherited disorders.

\medterm{Liver Damage} Liver damage is injury or loss of liver-cell function caused by infection, alcohol, medicines, toxins, fatty disease, autoimmune disease, or other processes.

\medterm{Liver Disease} Any condition damaging or impairing liver structure or function. Causes include infections, alcohol, medicines, toxins, fat buildup, immune disease, inherited disorders, blocked bile flow, and cancer.


\medterm{Liver Disease Diet} An individualized eating plan for a person with liver disease. It may address adequate energy and protein, sodium or fluid intake, alcohol avoidance, and nutrient deficiencies according to the condition.

\medterm{Liver Disorder} Any disease affecting the liver, including fatty liver, viral hepatitis, autoimmune disease, scarring, blockage of bile flow, or cancer. Symptoms may be absent or include jaundice, fatigue, pain, swelling, or abnormal blood tests.
\medterm{Liver Dysfunction} Liver dysfunction is reduced ability of the liver to perform tasks such as processing nutrients, making proteins, clearing toxins, and producing bile.

\medterm{Liver Enzymes} Blood measurements including ALT, AST, alkaline phosphatase, and GGT that can suggest liver or bile-duct injury. Abnormal results require interpretation with symptoms, medicines, imaging, and other tests.

\medterm{Liver Failure} A life-threatening loss of the liver's ability to make proteins, process nutrients, remove harmful substances, and regulate blood clotting. It can cause jaundice, bleeding, confusion, swelling, and multiple-organ failure.

\medterm{Liver Failure, Acute} Alternate index label; see \textit{Acute liver failure}.

\medterm{Liver Fibrosis} Excessive scar formation in the liver after ongoing injury or inflammation. It can progress to cirrhosis, reduce liver function, and result from viral hepatitis, alcohol, fatty liver, immune disease, or other causes.

\medterm{Liver Fluke} A parasitic flatworm that infects the liver or bile ducts, including Clonorchis, Opisthorchis, and Fasciola species.
\medterm{Liver Function (Liver (Anatomy And Function))} Alternate index label for Liver (Anatomy and Function).

\medterm{Liver Function Tests} Blood tests used to assess liver injury, bile-flow abnormalities, and liver synthetic function. They commonly
include aminotransferases, alkaline phosphatase, bilirubin, albumin, and clotting tests;
results must be interpreted together with the clinical context.

\medterm{Liver Function Tests (LFTs)} Blood tests assessing liver-cell injury, bile flow, and synthetic function, commonly including aminotransferases, alkaline phosphatase, bilirubin, albumin, and clotting measures.
\medterm{Liver Health} The ability of the liver to perform its many jobs, including processing nutrients and medicines, making important blood proteins, producing bile, storing energy, and helping remove harmful substances. Alcohol, infections, obesity, medicines, and other diseases can damage the liver.

\medterm{Liver Infection} Infection of liver tissue, such as a viral hepatitis infection or a bacterial liver abscess. Symptoms may include fever, tiredness, nausea, right-upper-abdominal pain, jaundice, or confusion and require medical evaluation.
\medterm{Liver Lobule} A small repeating unit of liver tissue in which blood from branches of the portal vein and hepatic artery flows past liver cells toward a central vein, while bile flows in the opposite direction.

\medterm{Liver Mass} An unusual area found in the liver, often on an ultrasound, CT scan, or MRI. It may be a harmless cyst or growth, a primary liver cancer, or cancer that spread from elsewhere; further testing may be needed.

\medterm{Liver Metastases} Liver metastases are cancer deposits that spread to the liver from another primary tumor and may cause no symptoms or produce pain, weight loss, or liver dysfunction.

\medterm{Liver Neoplasm} A liver neoplasm is an abnormal growth in the liver that may be benign or malignant and may arise from liver cells, bile ducts, or metastasis from another cancer.

\medterm{Liver Pain} Discomfort in the upper-right abdomen that may arise from the liver, gallbladder, bile ducts, stomach, muscles, or other organs. Persistent pain, jaundice, fever, vomiting, or confusion requires medical assessment.
\medterm{Liver Regeneration} The liver's ability to restore lost tissue by growing and reorganizing remaining cells after injury or partial surgical removal.
\medterm{Liver Scan} A nuclear medicine or imaging examination of the liver used to evaluate its size, shape, tissue distribution, focal lesions, blood flow, or selected functional abnormalities.

\medterm{Liver Shunt} An abnormal or surgically created channel that diverts blood around the liver, such as a portosystemic shunt; it can reduce detoxification and contribute to hepatic encephalopathy.

\medterm{Liver Spot} A common lay term for a solar lentigo, a flat brown skin patch caused by chronic ultraviolet exposure.
\medterm{Liver Spots} Liver spots are flat brown areas of sun-related pigmentation, also called solar lentigines, and are not caused by liver disease.

\medterm{Liver Transplant} Surgical replacement of a diseased liver with a healthy liver or part of a liver from a donor.

\medterm{Liver Transplantation} Surgery to replace a severely diseased or failing liver with a healthy donor liver or part of one. It may be used for advanced liver disease, sudden liver failure, or selected liver cancers; lifelong medicines help prevent rejection.

\medterm{Liver, Diseases Of} Disorders of the liver, including hepatitis, fatty liver, cirrhosis, tumors, and biliary disease, that can impair metabolism and detoxification.
\medterm{Liver, Enlarged} Alternate index label; see \textit{Hepatomegaly}.

\medterm{Living Donor} Individual who donates an organ (kidney, liver segment, or other tissue) while alive for transplantation into a recipient.

\medterm{Living Will} A type of advance directive that records a person's preferences for future medical treatment if the person cannot communicate or decide.

\medterm{Living Will (Advance Directive)} Alternate index label; see \textit{Living Will}.

\medterm{Living Will (Advance Statements)} Alternate index label; see \textit{Living Will}.

\medterm{Living With A Chronic Illness -  Dealing With Feelings} Living with a chronic illness can bring grief, anxiety, anger, or low mood; support from healthcare professionals, counseling, and trusted people can help.

\medterm{Living With A Chronic Illness - Reaching Out To Others} Reaching out to family, friends, support groups, or healthcare professionals can reduce isolation and help with practical and emotional effects of chronic illness.

\medterm{Living With Endometriosis} Living with endometriosis may involve pain management, hormonal or surgical treatment, pacing activities, fertility counseling, and follow-up for changing symptoms.

\medterm{Living With Hearing Loss} Living with hearing loss may include hearing aids or other devices, communication strategies, hearing protection, and treatment of an underlying cause when possible.

\medterm{Living With Heart Disease And Angina} Living with heart disease and angina involves prescribed medicines, activity and risk-factor management, follow-up, and urgent care for new or worsening chest pain.

\medterm{Living With Uterine Fibroids} Living with uterine fibroids may involve monitoring, medicines for bleeding or pain, procedures, or surgery depending on symptoms, size, location, and fertility plans.

\medterm{Living With Vision Loss} Living with vision loss may involve treating its cause, low-vision rehabilitation, assistive devices, home safety changes, and support with daily activities.


\medterm{Livor Mortis (Lividity)} Post-mortem purplish-red discoloration of dependent body parts resulting from
gravitational settling of blood in uncompressed capillaries and small veins following
cessation of circulation.

\synonyms
Lividity

\medterm{Livor Mortis (Postmortem Lividity)} Purplish-red discoloration of dependent body parts following death, caused by
gravitational pooling of blood in uncompressed capillaries and venules. Starts within 30
min--2 hours, becomes fixed after 6--12 hours.

\synonyms
Postmortem Lividity

\medterm{Lizard Bite} A puncture or tissue injury caused by a lizard; care includes cleaning and assessing for infection, tetanus, venom exposure, and significant tissue damage according to the species and circumstances.

\medterm{LLETZ} A procedure that uses a thin electrically heated wire loop to remove abnormal tissue from the cervix. The tissue is examined in the laboratory and the procedure can treat precancerous cervical changes.

\synonyms
LEEP (Loop Electrosurgical Excision Procedure)

\medterm{Lloyd-Davies Position} A modified lithotomy position in which the hips and knees are flexed and the legs are supported, commonly for pelvic or colorectal surgery.

\medterm{Loa Loa} The African eye worm, a filarial nematode transmitted by Chrysops deer flies, endemic to
the rainforests of West and Central Africa. Adult worms migrate subcutaneously, causing
Calabar swellings, and maycross the subconjunctiva of the eye, causing visible worm
migration. Microfilariae are detected in daytime peripheral blood.

\medterm{Loading Dose} The first, sometimes larger, dose of a medicine given to reach a helpful amount in the body quickly. It is followed by regular maintenance doses and is selected according to the medicine and the patient's condition.

\medterm{Lobaplatin} Lobaplatin is a platinum-based antineoplastic drug used in some countries and studied for its ability to damage cancer-cell DNA.

\medterm{Lobar Hemorrhage} Bleeding into one of the large regions, or lobes, of the brain. It can result from long-term high blood pressure, fragile blood vessels, blood-thinning medicines, injury, or other conditions and may cause sudden headache, weakness, confusion, seizures, or loss of consciousness.

\medterm{Lobar Pneumonia} Pneumonia affecting a large section or an entire lobe of the lung. It can cause fever, cough, chest pain, breathlessness, and low oxygen and may be caused by bacteria, viruses, or other organisms.

\medterm{Lobe} A lobe is a distinct subdivision of an organ or structure, such as a lung, liver, brain, thyroid, or gland.

\medterm{Lobectomy} Surgical removal of a lobe of an organ. Most commonly lung lobectomy (removal of one
lobe of the lung, performed for lung cancer, bronchiectasis, abscess, or benign tumor)
or temporal lobectomy (removal of the anterior temporal lobe for drug-resistant temporal
lobe epilepsy).

\medterm{Lobotomy} A historical brain operation that cut or damaged nerve connections in the front part of the brain. It was once used for severe mental illness but caused serious changes in personality and ability and is no longer an accepted treatment.


\medterm{Lobular Carcinoma} Lobular carcinoma is a cancer arising from cells of a breast lobule, the gland that produces milk; it may be invasive or confined to the lobule.

\medterm{Lobular Carcinoma In Situ} Abnormal cells confined to the breast lobules that increase the risk of developing invasive breast cancer.

\medterm{Lobule} A lobule is a small subdivision of an organ or tissue, such as a liver lobule, lung lobule, or breast lobule.

\medterm{Lobules} Small sections that make up an organ. In the breast, lobules are groups of glands that can make milk; in other organs, the meaning depends on the organ's structure and function.

\medterm{Local} Limited to a particular body site or region rather than involving the whole body. A local treatment or symptom affects the treated or affected area, although nearby or wider effects can sometimes occur.
\medterm{Local Anaesthesia} Loss of sensation in a limited body area caused by a local anesthetic without loss of consciousness.
\medterm{Local Anesthesia} Temporary loss of feeling in a limited area after a local anesthetic is applied or injected, without causing unconsciousness. Too much medicine or accidental injection into a blood vessel can cause seizures, heart problems, or other toxicity.

\medterm{Local Anesthetic Systemic Toxicity} A dangerous reaction when a local numbing medicine enters the bloodstream in excessive amounts. Early signs include tingling around the mouth, ringing in the ears, dizziness, or agitation; seizures and heart problems can follow.

\medterm{Local Anesthetic Toxicity} Harm from too much local anesthetic in the bloodstream. Early warning signs include ringing in the ears, numbness around the mouth, dizziness, or agitation; seizures, low blood pressure, abnormal rhythms, or cardiac arrest can follow and require emergency treatment.

\medterm{Local Invasion} Growth of a tumor into nearby tissues rather than only at its original site. It is a feature of malignant behavior and differs from distant metastasis.

\medterm{Local Therapy} Local therapy treats a disease at or near the affected area rather than throughout the body, for example with a cream, injection, surgery, or focused radiation.

\medterm{Local Treatment} Treatment directed at a specific body site or lesion rather than the whole body, such as a topical medicine, local injection, radiation field, or localized procedure.

\medterm{Localised} Confined to a limited area or a small number of adjacent body sites rather than widespread.

\medterm{Localization} Determining the body site responsible for a symptom, sign, lesion, or function. Accurate localization helps guide examination, imaging, diagnosis, treatment, surgery, rehabilitation, and assessment of nerve or organ damage.
\medterm{Localized} Limited to one particular area, organ, or body site without evidence of distant spread. A localized disease may still be serious, but treatment can sometimes focus on the affected area.

\synonyms
Localized disease

\medterm{Localized Edema} Swelling confined to a particular body region because of fluid accumulation, commonly caused by venous or lymphatic obstruction, inflammation, trauma, allergy, or local infection.

\medterm{Localized Scleroderma} An autoimmune skin disorder causing localized thickening and hardening of the skin, with or without involvement of deeper tissues; it differs from systemic sclerosis.

\medterm{Lochia} The normal vaginal discharge after childbirth. It is usually red and bloody at first, then becomes pink or brown, and finally yellowish or white over the following weeks. Heavy bleeding, large clots, fever, or foul odor needs assessment.

\medterm{Locked In Syndrome (Locked-In Syndrome)} Alternate index label for Locked-in Syndrome.

\medterm{Locked-In Syndrome} A condition in which a person remains conscious and able to think but is almost completely paralyzed and cannot speak. Some people can communicate by blinking or moving their eyes.

\synonyms
Locked In Syndrome (Locked-In Syndrome)

\medterm{Lockjaw} Difficulty opening the mouth because the jaw muscles are tight or painful. It may occur with tetanus, dental infection, jaw injury, or other conditions affecting the jaw muscles or joints.

\medterm{Lockjaw (Tetanus)} A serious neurologic illness caused by tetanospasmin produced by Clostridium tetani. It causes painful muscle stiffness and spasms, often beginning in the jaw.

\synonyms
Tetanus

\medterm{Locomotion} Movement of an organism or person from one place to another using coordinated muscles, bones, nerves, and balance systems.

\medterm{Locomotive System} The musculoskeletal and nervous systems working together to produce posture, balance, and movement from place to place.

\medterm{Locomotor Ataxia} An older term for tabes dorsalis, a neurologic complication of untreated syphilis causing sensory ataxia and lightning pains.
\medterm{Locomotor Ataxia (Tabes Dorsalis)} A late manifestation of untreated tertiary syphilis (neurosyphilis) caused by infection
with Treponema pallidum, involving progressive degeneration of the dorsal columns
(fasciculus gracilis and cuneatus) and dorsal roots of the spinal cord.

\synonyms
Tabes Dorsalis

\medterm{Locoregional} Affecting a defined local area and nearby tissues or lymph nodes, without necessarily involving the whole body.

\medterm{Locus} The specific location of a gene or other DNA sequence on a chromosome. It is described using the chromosome number and a coded position on its short or long arm.

\medterm{Locus Ceruleus} A small brainstem region containing nerve cells that release norepinephrine. It helps regulate alertness, attention, stress responses, sleep, and some automatic body functions.

\synonyms
Locus coeruleus

\medterm{Locus Coeruleus} The locus coeruleus is a brainstem nucleus whose noradrenaline-producing neurons influence arousal, attention, stress, sleep, and pain.

\medterm{Locus Minoris Resistentiae} Latin for a site of lesser resistance, meaning a body area considered more vulnerable to injury, infection, or disease because of local weakness or prior damage.

\medterm{Lod Score} A statistical measure used in genetic linkage analysis to compare how likely it is that two genetic markers are inherited together rather than independently. Higher positive scores provide stronger evidence of linkage in the studied family.

\medterm{Loeys-Dietz Syndrome} Loeys-Dietz syndrome is an inherited connective-tissue disorder associated with arterial aneurysm or dissection, distinctive skeletal and craniofacial features, allergic disease, and variable organ involvement. Lifelong vascular imaging and coordinated specialist care are important.

\synonyms
Syndrome Loeys-Dietz (Loeys-Dietz Syndrome)

\medterm{Loiasis} Loiasis is infection with the filarial worm Loa loa, transmitted by deerflies and sometimes causing transient eye-worm migration and localized swelling.

\medterm{Loin} The area of the lower back between the ribs and pelvis, including the side of the body called the flank. Pain there may come from muscles, the spine, kidneys, or nearby organs.

\medterm{Lomotil Overdose} Taking too much diphenoxylate–atropine can cause extreme sleepiness, confusion, dry mouth, fast heartbeat, slowed breathing, coma, or delayed worsening. It requires immediate emergency or poison-control advice, especially in children.

\medterm{Lomustine} Lomustine is an oral chemotherapy medicine used for selected brain tumors and other cancers; delayed low blood counts are an important risk.

\medterm{Lonafarnib} A farnesyltransferase-inhibiting medicine used for selected progeroid syndromes, in which abnormal protein processing contributes to premature aging features. It can cause gastrointestinal symptoms, liver abnormalities, and drug interactions.
\medterm{Loneliness} Subjective feeling of isolation. Associated with depression, cognitive decline,
mortality.

\medterm{Long Arm Of A Chromosome} The longer segment of a chromosome, designated q, extending from the centromere toward the chromosome's end.

\medterm{Long Bones} Bones that are longer than they are wide, such as the femur, tibia, humerus, and radius. Their shafts provide support and leverage for movement, while their ends form or help form joints; bone marrow inside them produces blood cells.

\medterm{Long COVID} Ongoing or new health problems that begin or continue after COVID-19. Symptoms may include tiredness, worsening after activity, difficulty thinking, breathlessness, changes in smell or taste, sleep problems, or effects on several organs.

\medterm{Long QT Syndrome} A heart rhythm disorder in which the heart takes too long to reset between beats. It may be inherited or caused by medicines or abnormal body salts and can lead to fainting, dangerous rhythms, or sudden death.

\medterm{Long-Acting Beta-2 Agonists} Inhaled medicines that keep the airways relaxed for many hours and reduce breathing symptoms in asthma or chronic obstructive pulmonary disease. For asthma, they are used with an inhaled corticosteroid, not alone.

\medterm{Long-Sightedness} Hyperopia, a refractive error in which distant objects may be clearer than near objects and accommodation is required for focusing.
\medterm{Long-Term Care} Ongoing help with daily activities, nursing, medical needs, or supervision for a person who cannot live fully independently. Care may be provided at home, in assisted living, or in a nursing facility and should match changing needs.

\medterm{Long-Term Care Facility} A residential or care setting providing ongoing assistance, nursing, rehabilitation, or personal care to people who cannot live independently or need sustained support.

\medterm{Long-Term Complications Of Diabetes} Chronic complications of diabetes caused mainly by vascular and metabolic injury, including retinopathy, kidney disease, neuropathy, cardiovascular disease, and foot disease.

\medterm{Long-Term Effect} A health effect that develops or persists after prolonged exposure, treatment, or disease. It may be beneficial, harmful, or uncertain.

\medterm{Long-Term Memory} Information retained for a long period, from hours or days to a lifetime. It includes explicit memories of facts and events and implicit skills or habits.

\medterm{Long-Term Potentiation} Long-term potentiation is a lasting increase in the strength of communication between nerve cells after repeated stimulation and is important in learning and memory.

\medterm{Long-Term Rehabilitation} Ongoing rehabilitation for chronic disease, lasting disability, or persistent injury. Goals include maintaining function, preventing complications, adapting activities, and supporting independence and quality of life.


\medterm{Longitudinal Section} A cut or image made along the long axis of an organ, limb, vessel, or other structure.

\medterm{Longitudinal Studies} Longitudinal studies follow participants or measurements over time to examine change, incidence, exposures, outcomes, and temporal relationships.

\medterm{Longitudinal Vaginal Septum} A wall of tissue that divides the vagina lengthwise, present from birth. It may cause painful sex, difficulty inserting tampons, blocked menstrual flow, or no symptoms. It can occur with an unusually shaped uterus and may be removed surgically when it causes problems.

\medterm{Loop Diuretic} A strong water tablet that makes the kidneys remove extra salt and water in urine. It is used for fluid buildup, heart failure, some kidney problems, and high blood pressure; dehydration and low body salts can occur.
\medterm{Loop Diuretics} A group of strong water tablets that make the kidneys remove salt and water quickly. They treat fluid buildup and some heart, kidney, or liver disorders; excessive use can cause dehydration, low salts, low blood pressure, or kidney problems.

\medterm{Loop Electrosurgical Excision Procedure} A procedure that uses a thin heated wire loop to remove abnormal tissue from the cervix. It both treats precancerous cervical changes and provides tissue for laboratory examination. Bleeding, infection, narrowing of the cervix, and a small increased risk of later premature birth can occur.

\synonyms
LEEP; Large Loop Excision of the Transformation Zone

\medterm{Loop Of Henle} A U-shaped part of each kidney’s tiny filtering tube. It helps the kidney conserve water and control the concentration of urine by allowing water and salts to move differently along its two limbs.

\medterm{Loose Bodies} Free fragments of cartilage, bone, or other tissue within a joint or body cavity that may cause pain, catching, or locking.
\medterm{Loperamide} Loperamide is a peripherally acting opioid-receptor agonist that slows intestinal motility and reduces diarrhoea; it should be avoided in bloody or febrile diarrhoea and excessive doses can cause severe cardiac arrhythmias.

\medterm{Loperamide Hydrochloride} Loperamide hydrochloride is an antidiarrheal medicine that slows movement through the intestine; excessive doses can cause dangerous heart rhythms.

\medterm{Lopinavir} Lopinavir is an HIV protease inhibitor used with ritonavir in antiretroviral treatment; drug interactions are common.

\medterm{Loratadine} An antihistamine used to relieve sneezing, runny nose, itchy or watery eyes, and hives. It usually causes less drowsiness than older antihistamines, but sleepiness, headache, or dry mouth can still occur.

\medterm{Lorazepam} Lorazepam is a benzodiazepine used short term for severe anxiety, seizures, or sedation; it can cause drowsiness, dependence, and breathing suppression.

\medterm{Lordosis} The inward curve of the neck or lower back. A curve that is unusually pronounced may be related to posture, pregnancy, muscle imbalance, spinal changes, or another condition and may cause back discomfort.

\medterm{Lordosis - Lumbar} Lumbar lordosis is the inward curve of the lower spine; excessive or reduced curvature may be postural, structural, degenerative, or related to pain.

\medterm{Losartan} Losartan is an angiotensin-receptor blocker used for high blood pressure, some kidney disease, and heart failure; it can raise potassium levels.

\medterm{Losartan Potassium} Losartan potassium is the salt form of losartan, an angiotensin-receptor blocker used to lower blood pressure and protect selected patients with kidney or heart disease.



\medterm{Loss Of Appetite} Reduced desire to eat. It may result from illness, medicines, mood disorders, pain, or other causes.

\medterm{Loss Of Bladder Control} Alternate index label; see \textit{Urinary incontinence}.

\medterm{Loss Of Brain Function - Liver Disease} Hepatic encephalopathy, a brain dysfunction caused by severe liver disease or portosystemic shunting and characterized by changes in attention, behavior, consciousness, or neuromuscular function.

\medterm{Loss Of Heterozygosity} Loss of heterozygosity is loss of one of two different versions of a gene in a cell, a change that can contribute to cancer when a tumor-suppressor gene is affected.

\medterm{Loss Of Sense Of Smell And Taste With COVID-19} COVID-19 can temporarily impair smell and taste through infection or inflammation of nasal and olfactory pathways. Most people improve, but persistent loss, distortion, poor nutrition, or neurologic symptoms warrants follow-up and smell-retraining advice.

\medterm{Loss Of Smell (Anosmia)} Partial or complete inability to detect odors, commonly caused by nasal obstruction, infection, inflammation, head injury, neurologic disease, or certain medicines.

\medterm{Loss-Of-Function Mutation} A genetic change that reduces or abolishes the normal activity or production of a gene product; its clinical effect depends on the gene, variant, and inheritance pattern.

\medterm{Loss-Of-Function Variant} A genetic variant that reduces or abolishes the normal activity of the encoded gene
product. Disease risk depends on the gene, inheritance pattern, residual function, and
whether one or both copies are affected.

\medterm{Loteprednol Etabonate} Loteprednol etabonate is a corticosteroid eye medicine used for eye inflammation; prolonged use can raise eye pressure and increase infection risk.

\medterm{Lou Gehrig Disease} Amyotrophic lateral sclerosis, a progressive neurodegenerative disease affecting upper and lower motor neurons and causing weakness, muscle wasting, and eventual respiratory impairment.

\medterm{Lou Gehrig'S Disease} Alternate index label; see \textit{Amyotrophic lateral sclerosis}.

\medterm{Louse} A small wingless insect that lives on hair, skin, or clothing and feeds on blood. Human lice include head, body, and pubic lice. They cause itching and may spread through close contact; treatment depends on the type and includes removing live lice and their eggs.

\medterm{Lovastatin} Lovastatin is a statin medicine that lowers LDL cholesterol and reduces cardiovascular risk; muscle injury and liver effects are uncommon but important risks.

\medterm{Low Back Pain} Pain or discomfort in the lower back, commonly caused by muscle or ligament strain, disc disease, arthritis, or spinal
nerve irritation. Weakness, numbness, new bladder or bowel problems, fever, major trauma, or unexplained
weight loss are warning signs requiring prompt medical assessment.

\synonyms
Lumbago (Low Back Pain)

\medterm{Low Back Pain - Acute} Pain in the lower back lasting less than about four weeks, commonly from muscle or ligament strain. Numbness, weakness, fever, injury, cancer history, or bladder problems require urgent assessment.

\medterm{Low Back Pain - Chronic} Lower-back pain lasting three months or longer. Causes include persistent strain, arthritis, disc disease, nerve irritation, inflammatory disease, or less commonly cancer; treatment addresses the cause, function, activity, and pain.

\medterm{Low Birth Weight} A birth weight below 2,500 grams (about 5 pounds 8 ounces), whether caused by being born early, slower growth before birth, or both. Lower weights carry greater risks and need careful newborn monitoring.

\medterm{Low Blood Oxygen (Hypoxemia)} An abnormally low oxygen level in arterial blood, usually related to lung disease, impaired breathing, circulation problems, or reduced oxygen availability.

\medterm{Low Blood Potassium} Low blood potassium, or hypokalemia, can cause weakness, cramps, constipation, abnormal heart rhythms, or paralysis and may result from losses, medicines, or hormonal disease.

\medterm{Low Blood Pressure} Low blood pressure can cause dizziness, fainting, weakness, or shock and may result from dehydration, medicines, bleeding, infection, heart disease, or autonomic dysfunction.

\medterm{Low Blood Pressure (Hypotension) Causes} Alternate index label; see \textit{Hypotension}.

\medterm{Low Blood Sodium} Low blood sodium, or hyponatremia, reflects excess water relative to sodium and can cause headache, confusion, seizures, or coma when severe or rapidly developing.

\medterm{Low Blood Sugar} Low blood sugar, or hypoglycemia, occurs when circulating glucose is insufficient for normal brain and body function, causing sweating, tremor, confusion, seizures, or loss of consciousness.

\medterm{Low Blood Sugar - Newborns} Low blood sugar in a newborn can cause jitteriness, poor feeding, lethargy, apnea, or seizures and is more likely with prematurity, growth problems, maternal diabetes, or illness.


\medterm{Low Body Temperature} Alternate index label; see \textit{Hypothermia}.

\medterm{Low Calcium Level  - Infants} An abnormally low calcium level in a baby. It may cause jitteriness, poor feeding, sleepiness, pauses in breathing, seizures, or abnormal heart rhythms and requires age-appropriate medical assessment.

\medterm{Low FODMAP Diet} A temporary eating plan that limits certain poorly absorbed carbohydrates that can cause gas, bloating, pain, or diarrhea in some people with irritable bowel syndrome. Foods are gradually reintroduced with dietitian guidance to identify personal triggers.

\medterm{Low Grade} Describes a disease or tumor that appears less aggressive or more slowly developing than higher-grade forms. The exact meaning and outlook depend on the specific disease and test used.

\medterm{Low Hemoglobin Count} A hemoglobin level below the expected range, often indicating anemia from blood loss, reduced production, nutritional deficiency, inflammation, or increased destruction of red blood cells.

\medterm{Low Molecular Weight Heparin} A blood-thinning medicine used to prevent or treat blood clots. It is usually injected under the skin and can cause bleeding, bruising, or, rarely, a serious fall in platelets.

\medterm{Low Nasal Bridge} A low nasal bridge is reduced height of the nasal root and may be a normal familial feature or part of a congenital, genetic, or developmental syndrome.

\medterm{Low Neutrophil Count (Neutropenia)} Alternate index label for Neutropenia.

\medterm{Low Platelet Count} Alternate index label; see \textit{Thrombocytopenia}.

\medterm{Low Potassium (Hypokalemia)} An abnormally low potassium level in the blood, which can cause weakness, cramps, abnormal heart rhythms, or paralysis and may result from losses, medicines, or hormonal disorders.

\medterm{Low Sex Drive In Women} A lasting decrease in sexual interest or desire that causes personal distress or difficulty in a relationship. Possible causes include stress, depression, pain, relationship problems, hormonal changes, illness, or medicines; treatment depends on the cause and the person's wishes.

\medterm{Low Sperm Count} A semen condition in which sperm concentration is below the expected reference range. It may reduce the chance of pregnancy, but fertility also depends on sperm movement, shape, ejaculation, and the partner's health.

\medterm{Low Testosterone} Alternate index label; see \textit{Male hypogonadism}.

\medterm{Low Testosterone (Hypogonadism)} A condition in which a person has symptoms of testosterone deficiency and repeatedly low morning testosterone levels. Causes may involve the testes, ovaries, pituitary gland, medicines, obesity, or another illness and need medical evaluation.

\synonyms
Hypogonadism

\medterm{Low Testosterone (Low T)} Alternate index label for Low T.

\medterm{Low Vision} Visual impairment that cannot be fully corrected with spectacles, contact lenses,
medication, or surgery, impairing daily activities. Managed with low-vision aids
(magnifiers, telescopes, adaptive devices), lighting, and occupational therapy.

\medterm{Low White Blood Cell Count} A reduced number of white blood cells, which can increase susceptibility to infection and may result from medicines, infection, autoimmune disease, nutritional deficiency, or bone marrow disorders.

\medterm{Low White Blood Cell Count And Cancer} A reduced number of infection-fighting blood cells in a person with cancer or receiving cancer treatment. The risk of serious infection depends on which cells are low and how severely they are reduced.

\medterm{Low-Calorie Diet} An eating plan that provides less energy than a person usually uses, often to support weight loss or a medical goal. It should still provide adequate protein, vitamins, minerals, fluids, and fiber.
\synonyms
Calorie-restricted diet

\medterm{Low-Density Lipoprotein} A particle that carries cholesterol through the blood. Too much LDL cholesterol can collect in artery walls, forming fatty deposits that narrow arteries and increase the risk of heart attack and stroke.
\medterm{Low-Dose CT Scan} A CT scan performed with less radiation than a standard scan. It is used for purposes such as lung-cancer screening or follow-up of lung nodules, but abnormal findings may require additional tests.

\medterm{Low-Fiber Diet} An eating plan that limits foods such as whole grains, nuts, seeds, raw vegetables, and fruit skins. It may temporarily reduce stool volume for selected bowel problems or medical procedures.

\medterm{Low-Flow Low-Gradient Aortic Stenosis} A severe narrowing of the aortic valve in which less blood moves through the valve and the pressure difference is lower than expected. Echocardiography and other tests help confirm severity.
\synonyms
LF-LG Aortic Stenosis; Low-Gradient Severe Aortic Stenosis

\medterm{Low-Flow Priapism} Alternate index label; see \textit{Priapism}.

\medterm{Low-Grade Appendiceal Mucinous Neoplasm} A mucin-producing appendiceal neoplasm with low-grade cytologic atypia that may dilate the appendix and rupture, disseminating mucin into the peritoneum as pseudomyxoma peritonei.

\medterm{Low-Grade Central Osteosarcoma} A rare low-grade malignant osteogenic tumor arising within the medullary cavity of bone. It may resemble fibrous dysplasia or a benign fibro-osseous lesion and requires complete oncologic excision.

\medterm{Low-Grade Endometrial Stromal Sarcoma} A malignant uterine stromal tumor resembling proliferative-phase endometrial stroma, characterized by infiltrative myometrial or vascular growth and often driven by characteristic gene fusions.
\medterm{Low-Phosphorus Diet} A dietary pattern limiting phosphorus intake for selected patients with chronic kidney
disease or persistent hyperphosphatemia. It emphasizes avoidance of phosphate additives
while maintaining adequate protein and overall nutritional status.

\medterm{Low-Potassium Diet} A dietary pattern limiting potassium intake when clinically indicated, usually for
hyperkalemia or advanced kidney disease. It must be individualized because unnecessary
restriction can reduce diet quality and cause inadequate intake of fruits and
vegetables.

\medterm{Low-Residue Diet} A short-term diet limiting dietary fiber and foods that leave substantial undigested
material in the bowel. It may be used in selected gastrointestinal conditions or around
certain procedures but is not a universal treatment for inflammatory bowel disease.

\medterm{Low-Salt Diet} A low-salt diet limits sodium intake and may help control blood pressure or fluid retention in conditions such as heart failure and kidney disease.

\medterm{Low-Set Ear} An ear positioned lower on the side of the head than expected relative to facial landmarks; it may be an isolated variation or a feature of a genetic or developmental condition.

\medterm{Low-Set Ears And Pinna Abnormalities} Differences in the position, size, shape, or folds of the external ear, including ears positioned lower than usual. They may occur alone or with a congenital syndrome, so hearing and other physical findings may need assessment.

\medterm{Lower Airway Obstruction} Narrowing or blockage of the breathing tubes inside the lungs. Asthma, chronic lung disease, infection, swelling, or a foreign object may cause wheezing, cough, breathlessness, or reduced airflow.

\medterm{Lower Esophageal Ring} A thin ring of tissue near the junction of the esophagus and stomach, often called a Schatzki ring. It can narrow the esophagus and cause intermittent difficulty swallowing solid food or food becoming stuck.

\medterm{Lower Esophageal Sphincter} A ring of muscle where the esophagus joins the stomach. It opens during swallowing and normally closes afterward to limit backward flow of stomach contents and acid.
\synonyms
LES

\medterm{Lower Extremity} The lower extremity is the hip, thigh, knee, leg, ankle, and foot, including their bones, joints, muscles, nerves, vessels, and skin.

\medterm{Lower Extremity Venous Doppler} An ultrasound test of the leg veins that shows blood flow and whether the veins compress normally. It is commonly used to look for a deep-vein blood clot.

\medterm{Lower Leg} The part of the limb between the knee and ankle, containing the tibia, fibula, muscles, nerves, and blood vessels.

\medterm{Lower Motor Neuron Syndrome} A pattern of weakness caused by damage to nerves that carry movement signals from the spinal cord or brainstem to muscles. It may cause muscle wasting, twitching, low muscle tone, and reduced reflexes.

\medterm{Lower Segment Cesarian Section} A cesarean delivery performed through an incision in the lower uterine segment, usually through a low transverse uterine incision.

\medterm{Lower Urinary Tract Symptoms} Problems with storing or passing urine, including urgency, frequency, nighttime urination, weak stream, straining, leakage, or a feeling of incomplete emptying. Causes include infection, prostate enlargement, bladder disorders, nerve disease, medicines, and blockage.

\medterm{Lozenges} Solid medicinal preparations that dissolve slowly in the mouth, usually to provide local treatment for the throat or oral cavity.
\medterm{LSD} Lysergic acid diethylamide is a potent hallucinogen that can cause altered perception, visual hallucinations, autonomic effects, and unpredictable changes in mood and behavior.

\medterm{Lubiprostone} Lubiprostone activates intestinal chloride channels and is used for selected forms of chronic constipation and irritable bowel syndrome with constipation.

\medterm{Lubricant} A substance that reduces friction between surfaces; medical lubricants are used on skin, mucosa, instruments, or devices and should be selected for tissue compatibility and intended use.


\medterm{Lucilia} A genus of blowflies; some species can cause wound myiasis, while sterile larvae have occasionally been used in maggot debridement.
\medterm{Ludwig Angina} A rapidly spreading cellulitis of the floor of the mouth and submandibular spaces,
usually arising from an infected lower molar. Tongue elevation, dysphagia, drooling, and
airway obstruction can develop rapidly and require emergency management.

\medterm{Ludwig'S Angina} A rapidly progressive, life-threatening cellulitis of the submandibular, sublingual, and
submental spaces, typically originating from an odontogenic infection of the lower
second or third molars.


\medterm{Lugano Classification} A system used to describe how far Hodgkin or non-Hodgkin lymphoma has spread and how it responds to treatment. It combines examination, blood tests, and imaging, often PET with CT. Stages range from disease in one area to widespread disease; the results help guide treatment and follow-up.

\medterm{Lugol'S Iodine} An aqueous solution of iodine and potassium iodide used as an antiseptic, in selected thyroid treatments, and in some laboratory procedures.
\medterm{Lugol'S Solution} Lugol's solution is a mixture of iodine and potassium iodide used in selected laboratory, diagnostic, and medical applications; it can irritate tissue and affect thyroid function.

\medterm{Lumbago} A non-specific term for low back pain, typically referring to acute or chronic pain in
the lumbar region (between the lower ribs and the gluteal folds). Lumbago is a symptom,
not a diagnosis.

\medterm{Lumbago (Low Back Pain)} Alternate index label for Low Back Pain.

\medterm{Lumbar} Relating to the lower back, especially the region or five vertebrae between the thoracic spine and sacrum.

\medterm{Lumbar Disc Herniation} Protrusion or extrusion of lumbar disc material through the annulus fibrosus, which may irritate or
compress a nerve root and cause back or leg pain. Symptoms, examination findings, and
imaging determine its significance and treatment.

\medterm{Lumbar Lordosis Lower Spine} Lumbar lordosis is the normal inward curve of the lower spine; excessive or reduced curvature may be associated with posture, muscle imbalance, hip disease, vertebral deformity, or pain. Examination and imaging are guided by symptoms rather than appearance alone.

\medterm{Lumbar MRI Scan} Magnetic resonance imaging of the lumbar spine used to evaluate vertebrae, discs, spinal canal, nerve roots, marrow, and soft tissues in conditions such as radiculopathy or spinal stenosis.

\medterm{Lumbar Plexus} The lumbar plexus is a network of nerves from the upper lumbar spinal roots that supplies the lower abdomen, pelvis, and parts of the thigh and leg.

\medterm{Lumbar Puncture} A procedure that uses a needle in the lower back to collect cerebrospinal fluid, the liquid around the brain and spinal cord. The fluid can be tested for infection, bleeding, inflammation, or cancer; pressure can also be measured, and some medicines can be given. Clinicians check bleeding risk and possible increased pressure in the skull first. Headache, back pain, bleeding, and infection are possible complications.

\synonyms
LP
Spinal Tap

\medterm{Lumbar Radiculopathy} Irritation or pressure on a nerve leaving the lower spine, causing pain that travels into the buttock or leg. It may also cause numbness, tingling, or weakness and is often called sciatica.

\medterm{Lumbar Region} The lumbar region is the lower part of the back between the ribs and pelvis, containing the five lumbar vertebrae.

\medterm{Lumbar Spine} The lower part of the backbone between the ribs and pelvis, usually made of five vertebrae. It bears much of the body’s weight and contains nerves supplying the legs; discs, joints, or nerves there may cause lower-back pain or sciatica.

\synonyms
Lumbar region

\medterm{Lumbar Spine CT Scan} Computed tomography of the lumbar spine that provides detailed images of vertebral bone and alignment and is particularly useful for detecting fractures or other bony abnormalities.

\medterm{Lumbar Stenosis} Narrowing of the spinal canal or nerve openings in the lower back. It can squeeze spinal nerves and cause back pain, leg pain, numbness, weakness, or symptoms that worsen with walking.

\synonyms
Stenosis Lumbar (Lumbar Stenosis)
Stenosis Spinal (Lumbar Stenosis)

\medterm{Lumbar Vertebrae} The five large bones of the lower spine between the ribs and pelvis. They bear much of the body’s weight, allow bending, protect spinal nerves, and can be affected by fractures, arthritis, or disc problems.

\medterm{Lumbar Vertebrae (L1-L5)} The five lumbar vertebrae in the lower back, numbered L1 through L5, between the thoracic spine and sacrum and bearing much of the body's load.
\medterm{Lumbosacral Region} The lumbosacral region is the junction between the lower lumbar spine and the sacrum at the back of the pelvis.

\medterm{Lumbosacral Spine CT} Computed tomography of the lower spine and sacrum using x-rays and computer reconstruction to evaluate fractures, bone lesions, alignment, and selected causes of nerve compression.

\medterm{Lumbosacral Spine X-Ray} Plain radiographic imaging of the lumbar spine and sacrum used to evaluate alignment, fractures, degenerative changes, deformity, and some bone lesions.

\medterm{Lumbricals} Small hand or foot muscles arising from tendons of the flexor digitorum profundus or flexor digitorum longus and assisting flexion-extension coordination.
\medterm{Lumen} The inner space of a hollow organ, vessel, or tube.

\medterm{Luminescence} Luminescence is emission of light from a substance through a process other than high-temperature incandescence, such as fluorescence or chemiluminescence.

\medterm{Luminescence Type} Luminescence type identifies the mechanism by which a material emits light, such as fluorescence, phosphorescence, bioluminescence, or chemiluminescence.

\medterm{Lump In The Abdomen} An abdominal lump may be a hernia, enlarged organ, cyst, stool, inflammation, or tumor; pain, vomiting, fever, growth, or a pulsatile mass requires prompt assessment.

\medterm{Lumpectomy} Breast-conserving surgery that removes a cancer or other abnormal area with a small margin of nearby tissue. Radiation or additional treatment may be recommended afterward, depending on the diagnosis and risk of recurrence.

\medterm{Lumpy Jaw} A chronic granulomatous infection of the jaw, usually cervicofacial actinomycosis, causing swelling and sinus formation.
\medterm{Lunate} A small crescent-shaped bone in the center of the wrist. It helps the wrist move and can be injured by falls; loss of its blood supply can cause Kienböck disease and bone collapse.

\medterm{Lung} One of the two breathing organs in the chest. The lungs take oxygen into the blood and remove carbon dioxide. The right lung has three sections, and the left has two to leave room for the heart.

\medterm{Lung Abscess} A pocket of infection and pus that forms inside damaged lung tissue, often after inhaling stomach contents or because of bacteria. It may cause fever, cough, foul-smelling sputum, chest pain, and breathlessness.

\medterm{Lung Anatomy (Lungs Design And Purpose)} Alternate index label for Lungs Design And Purpose.

\medterm{Lung Cancer} Cancer that begins in the lungs or airways. The main types are small-cell and non-small-cell cancer. Smoking is the leading risk factor, but radon, asbestos, air pollution, inherited risk, and other exposures can also contribute. Early disease may cause no symptoms.
\synonyms
Cancer, Lung
Cancer In Lung (Lung Cancer)

\medterm{Lung Cancer - Small Cell} A fast-growing lung cancer made of small neuroendocrine cells that often spreads early.

\medterm{Lung Collapse} Alternate index label; see \textit{Atelectasis}; see also \textit{Pneumothorax}.

\medterm{Lung Compliance} How easily the lungs stretch when a person breathes in. Stiff lungs require more effort, while overly flexible lungs may not spring back well, as can occur in emphysema.

\medterm{Lung Cyst} A lung cyst is a round air- or fluid-containing space in lung tissue, caused by developmental disease, infection, emphysema, inflammation, or another disorder.

\medterm{Lung Development} Lung development is embryonic and postnatal formation and maturation of airways, alveoli, blood vessels, and gas-exchange capacity.

\medterm{Lung Diffusion Testing} A breathing test that measures how well oxygen-like gas moves from the lung air sacs into the blood. It helps assess diseases such as emphysema, lung scarring, or disorders of the lung blood vessels.

\medterm{Lung Disease} Lung disease includes disorders affecting the breathing tubes, air sacs, lung tissue, covering, blood vessels, or chest wall. These conditions may cause cough, breathlessness, chest discomfort, wheezing, or low oxygen.


\medterm{Lung Disease, Interstitial} Alternate index label; see \textit{Interstitial lung disease}.

\medterm{Lung Fibroma} A usually noncancerous growth made of fibrous tissue in or near the lung. It often causes no symptoms and is found on imaging, but its appearance, growth, and location may require further imaging or biopsy.

\medterm{Lung Fluke} A parasitic flatworm, usually of the Paragonimus group, acquired by eating undercooked freshwater crustaceans. Infection may cause chronic cough, chest pain, blood in sputum, or lung lesions and can rarely affect the brain.
\medterm{Lung Function Tests} Physiologic tests that measure aspects of breathing, including airflow, lung volumes, and gas transfer, to assess respiratory function.


\medterm{Lung Herniation} Protrusion of lung tissue through a weak or torn area of the chest wall, usually after chest trauma or surgery. It may appear as a bulge that enlarges with coughing and can cause pain or breathing difficulty.

\medterm{Lung Injury} Damage to lung tissue or the pulmonary circulation caused by trauma, infection, inhaled substances, inflammation, aspiration, oxygen toxicity, or other disease.

\medterm{Lung Lavage} A procedure in which doctors wash parts of the airways or lungs with sterile fluid and then remove it. It may clear certain material from the lungs or collect a sample for testing.

\medterm{Lung Metastases} Lung metastases are cancer deposits that have spread to the lungs from another primary tumor and may cause cough, breathlessness, chest pain, or no symptoms.

\medterm{Lung Needle Biopsy} CT- or ultrasound-guided percutaneous sampling of a lung or pleural lesion to evaluate selected masses, infection, or diffuse disease. Pneumothorax and pulmonary hemorrhage are important complications.

\medterm{Lung Neoplasm} An abnormal growth in the lung or airways that may be noncancerous or cancerous. It may cause cough, blood in sputum, chest pain, breathlessness, or no symptoms; imaging and sometimes tissue sampling determine its nature.

\medterm{Lung Nodule} A rounded, usually well-defined opacity in the lung measuring up to 3 cm on imaging. It may be benign, infectious, inflammatory, or an early manifestation of lung cancer and often requires follow-up assessment.

\medterm{Lung PET Scan} Positron-emission tomography, often combined with computed tomography, using a radiotracer to assess metabolic activity in lung nodules, cancer, infection, or inflammation.

\medterm{Lung Plethysmography} A breathing test performed in a sealed room to measure how much air the lungs hold, including air that cannot be breathed out, and how much resistance the airways provide.

\medterm{Lung Problems And Volcanic Smog} Volcanic smog contains gases and fine particles that irritate eyes and airways and can worsen asthma or lung disease; limiting exposure and following public-health advice reduce harm.

\medterm{Lung Protective Ventilation} A ventilator approach that uses gentle breath sizes and carefully controlled pressures to reduce further lung injury. It is especially important in acute respiratory distress and requires adjustment to oxygen levels, blood gases, and the person’s condition.

\medterm{Lung Surgery} An operation on the lungs or nearby structures to remove disease, repair injury, obtain tissue, or improve breathing. Procedures range from removing a small area to an entire lung; bleeding, infection, air leakage, and breathing problems are possible.


\medterm{Lung Transplant} Surgery to replace one or both severely diseased lungs with donor lungs. It may be considered for advanced lung failure from conditions such as emphysema, lung scarring, cystic fibrosis, or pulmonary hypertension and requires lifelong immune-suppressing medicines.

\medterm{Lung Transplantation} Surgical replacement of a diseased lung with a healthy lung from a donor. Indications
include IPF, COPD, CF, PAH, and bronchiectasis. Options include single lung transplant,
double lung transplant, or heart-lung transplant. Post-transplant immunosuppression is
required lifelong. Complications include rejection and infection.

\medterm{Lung, Diseases Of} Conditions affecting the airways, alveoli, pleura, or pulmonary circulation, including infection, asthma, COPD, fibrosis, and cancer.
\medterm{Lung Volume Reduction Surgery} An operation that removes the most damaged, overinflated parts of the lungs in carefully selected people with severe emphysema. This gives healthier lung tissue and the diaphragm more room to work, which may improve breathing and activity. It has important risks, including air leakage, infection, and respiratory failure.

\synonyms
LVRS; Reduction Pneumoplasty; Bilateral Lung Volume Reduction

\medterm{Lungs} The paired respiratory organs in the chest that exchange oxygen and carbon dioxide between air and blood; they also contribute to acid--base regulation.

\medterm{Lungs Design And Purpose} The lungs exchange oxygen and carbon dioxide through branching airways and microscopic alveoli surrounded by capillaries. The airways filter, warm, and humidify air, while respiratory muscles drive ventilation; disease can impair airflow, diffusion, or blood flow.

\synonyms
Lung Anatomy (Lungs Design And Purpose)

\medterm{Lungs Fluid (Pleural Effusion (Fluid In The Chest Or On Lung))} Alternate index label for Pleural Effusion (Fluid In the Chest or On Lung).

\medterm{Lunula} The pale crescent-shaped area at the base of a fingernail or toenail. It is the visible front part of the nail-forming region, although its size and visibility vary normally.
\medterm{Lupus} A group of autoimmune diseases in which the immune system attacks the body's own
tissues. Systemic lupus erythematosus may affect the skin, joints, kidneys, blood cells,
nervous system, and other organs. Also called SLE.

\synonyms
SLE.

\medterm{Lupus (Systemic Lupus)} Alternate index label for Systemic Lupus.



\medterm{Lupus Erythematosus} A long-term autoimmune disease in which the immune system attacks the body’s own tissues. It can cause tiredness, joint pain, rashes, and inflammation or damage in organs such as the kidneys, lungs, heart, brain, or blood vessels.

\medterm{Lupus Erythematosus (Systemic Lupus Erythematosus - SLE)} A chronic autoimmune disease that can affect many parts of the body, including the skin, joints, kidneys, heart, lungs, nerves, and blood cells. Symptoms often flare and improve, and treatment depends on the organs involved.

\synonyms
Systemic Lupus Erythematosus - SLE

\medterm{Lupus Nephritis} Kidney inflammation caused by systemic lupus erythematosus. It can cause blood or protein in the urine,
high blood pressure, and reduced kidney function. Classification is based largely on
kidney-biopsy findings, and treatment depends on the class, severity, kidney function,
and individual risks.

\medterm{Lupus Panniculitis} A chronic inflammatory form of cutaneous lupus affecting subcutaneous fat, producing firm deep nodules that may heal with lipoatrophy, scarring, or overlying discoid changes.

\medterm{Lupus Pernio} Persistent red-purple or blue-purple swelling and discoloration of the nose, cheeks, lips, ears, or fingers, usually associated with sarcoidosis. It often indicates longer-lasting disease and may cause scarring or cosmetic changes.

\medterm{Lupus Vulgaris} A long-lasting form of tuberculosis affecting the skin. It usually causes slowly enlarging reddish-brown patches or plaques that may scar and is treated with a combination of tuberculosis medicines.


\medterm{Luteal Cell} A cell in the ovarian corpus luteum that produces progesterone after ovulation and supports the uterine lining.

\medterm{Luteal Phase Defect} A term used when the part of the menstrual cycle after ovulation may not provide enough progesterone or enough time for the uterine lining to support pregnancy. The diagnosis is uncertain and should not be made from cycle length alone.

\synonyms
Luteal Phase Deficiency; LPD

\medterm{Lutein} A yellow plant pigment found in leafy green vegetables and concentrated in the retina. Along with related pigments, it helps filter blue light and may protect eye tissues from oxidative damage.

\medterm{Luteinization} Luteinization is the transformation of an ovarian follicle after ovulation into a corpus luteum that produces progesterone.

\medterm{Luteinizing Hormone} A hormone released by the pituitary gland. In people with ovaries, a mid-cycle rise triggers ovulation and supports hormone production; in people with testes, it stimulates testosterone production. Blood testing helps assess puberty, fertility, and pituitary or reproductive disorders.

\medterm{Luteinizing Hormone Blood Test} A blood test measuring luteinizing hormone, which stimulates ovulation and ovarian hormone production in females and testosterone production in males; it helps evaluate puberty, fertility, and pituitary or gonadal disorders.

\medterm{Luteinizing Hormone-Releasing Hormone} An older name for gonadotropin-releasing hormone (GnRH), a hypothalamic hormone released in pulses that stimulates the anterior pituitary to secrete luteinizing hormone and follicle-stimulating hormone.

\medterm{Luteolysis} Luteolysis is the breakdown of the corpus luteum, causing progesterone levels to fall when pregnancy has not occurred.

\medterm{Luteoma} A luteoma is a usually benign ovarian mass of hormone-producing lutein cells, sometimes associated with pregnancy and androgen excess.

\medterm{Luteoma Of Pregnancy} A benign, pregnancy-associated ovarian proliferation of luteinized stromal cells driven by hormonal stimulation. It may cause maternal or fetal androgen excess and usually regresses after delivery.

\medterm{Luteotropic Hormone} A hormone that supports the corpus luteum; the term traditionally refers mainly to prolactin, although usage varies by species and source.
\medterm{Lutetium-177} A radioactive form of lutetium that releases therapeutic radiation and a small amount of detectable radiation. Attached to a targeting molecule, it can deliver treatment to selected neuroendocrine or prostate cancer cells while limiting exposure to other tissues.

\medterm{Luxation} Complete dislocation of a joint, where the bone ends lose normal contact. It can stretch or tear nearby ligaments, nerves, and blood vessels, causing severe pain, deformity, swelling, and inability to move.

\medterm{LVAD} Left-ventricular assist device, an implanted pump helping the heart’s main chamber circulate blood. It may support people with severe heart failure while awaiting transplant or as long-term treatment.

\synonyms
Left-ventricular assist device

\medterm{LVNC} Alternate name; see \textit{Left Ventricular Noncompaction}.

\medterm{LVRS} Alternate name; see \textit{Lung Volume Reduction Surgery}.

\medterm{Lyase} An enzyme that breaks or joins certain chemical bonds without adding water or using oxidation as the main step. Lyases participate in many processes that make or use energy in cells.


\medterm{Lyme Disease} A bacterial infection spread by infected ticks; blacklegged ticks transmit it in the United States. Early symptoms may include an expanding rash, fever, tiredness, headache, and muscle or joint aches; untreated infection can later affect nerves, joints, or the heart.

\synonyms
Arthritis Lyme (Lyme Disease)


\medterm{Lyme Disease Blood Test} A blood test that looks for immune proteins made after infection with Lyme disease bacteria. Results depend on symptoms, tick exposure, and timing; early testing can be negative and a positive result does not always prove active infection.


\medterm{Lymph} A clear fluid collected from body tissues and carried through lymph vessels. It transports immune cells and excess fluid to lymph nodes and eventually returns it to the bloodstream.

\medterm{Lymph Leakage} Lymph leakage is escape of lymphatic fluid after injury, surgery, or obstruction and may cause swelling, fluid collections, or chylous drainage.

\medterm{Lymph Node} A lymph node is a small immune organ that filters lymph and helps activate immune responses to infection, inflammation, or abnormal cells.

\medterm{Lymph Node / Gland} A small immune organ along a lymph vessel that filters lymph and helps the body recognize infection, inflammation, or abnormal cells. Nodes are found in groups in the neck, armpits, chest, abdomen, and groin.

\medterm{Lymph Node Biopsy} Removal of part or all of a lymph node for microscopic and sometimes microbiologic or molecular examination to evaluate malignancy, infection, or inflammatory disease.

\medterm{Lymph Node Culture} Laboratory culture of lymph-node tissue or aspirate to identify bacteria, fungi, or mycobacteria causing lymphadenitis; culture is interpreted with histology and other tests.

\medterm{Lymph Node Dissection} Surgery that removes one or more groups of lymph nodes to check for cancer or treat cancer spread. Removed tissue is examined in a laboratory, and the operation can cause swelling or numbness.

\medterm{Lymph Node Structure} A lymph node is a small, bean-shaped immune organ with a capsule, spaces that carry lymph, and areas rich in immune cells. It filters lymph and may enlarge during infection, inflammation, or cancer spread.

\medterm{Lymph Nodes} Small immune organs filtering lymph fluid and helping detect infections or abnormal cells. They may enlarge or become tender with infection, inflammation, immune disease, or cancer spread.

\synonyms
Lymph glands

\medterm{Lymph Nodes, Swollen} Alternate index label; see \textit{Swollen lymph nodes}.

\medterm{Lymph Swollen Glands (Swollen Lymph Nodes)} Alternate index label for Swollen Lymph Nodes.

\medterm{Lymph Swollen Nodes (Swollen Lymph Nodes)} Alternate index label for Swollen Lymph Nodes.

\medterm{Lymph System} The lymph system includes lymph, lymphatic vessels, lymph nodes, spleen, thymus, and related tissues that return fluid and support immune defense.

\medterm{Lymphadenectomy} Surgical removal of lymph nodes, either sentinel or regional, for staging or therapeutic
purposes; examples include axillary, inguinal, cervical, and retroperitoneal lymph node
dissection.

\medterm{Lymphadenitis} Inflammation or infection of one or more lymph nodes, causing swelling and sometimes tenderness, warmth, or pus. Causes include viruses, bacteria, cat-scratch disease, tuberculosis, and other disorders.

\medterm{Lymphadenitis, Mesenteric} Alternate index label; see \textit{Mesenteric lymphadenitis}.

\medterm{Lymphadenopathy} Enlargement of one or more lymph nodes. It may be localized or widespread and can result from infection, inflammation, immune disease, medicines, or cancer; persistent, hard, or unexplained nodes need assessment.

\medterm{Lymphangiectasis} Abnormal widening of lymph vessels. It may be present from birth or develop later and can cause swelling, leakage of lymph, loss of protein through the intestine, or fluid around the lungs.

\medterm{Lymphangiogenesis} Formation of new lymph vessels from existing vessels. This can help repair tissue and drain fluid, but it may also contribute to chronic swelling, inflammation, or the spread of cancer.

\medterm{Lymphangiogram} A lymphangiogram uses contrast imaging to show lymphatic vessels and nodes, helping evaluate obstruction, leakage, malformation, or spread of disease.

\medterm{Lymphangiography} An imaging test that uses contrast dye to show lymphatic vessels and their flow. It can help find blocked vessels, leaks, or abnormal connections, although newer scans are often preferred.

\medterm{Lymphangioleiomyomatosis} A rare lung disease in which unusual cells grow in the lungs and lymphatic vessels, causing many small cysts. It mainly affects women and can lead to breathlessness, collapsed lung, bleeding, or kidney tumors.

\medterm{Lymphangioma} A noncancerous birth defect in which lymph vessels form abnormally and become enlarged. It usually appears in infancy or childhood as a soft swelling, often in the neck, armpit, or chest.

\medterm{Lymphangiomatosis} A rare disorder in which abnormal lymph vessels grow in one or more organs. It may cause swelling, fluid buildup, bleeding, bone damage, or breathing problems, depending on the tissues involved.

\medterm{Lymphangiomyoma} A rare, usually noncancerous growth made of smooth-muscle-like cells around lymph vessels. It is most often associated with lymphangioleiomyomatosis and may contribute to cysts or blockage of lymph flow.
\medterm{Lymphangiosarcoma} A rare, fast-growing cancer arising from the lining of lymph vessels. It may develop in tissue with long-standing lymph swelling and can appear as a bruise-like, purple, or ulcerated skin lesion.

\medterm{Lymphangitis} Inflammation and infection of lymphatic vessels, usually spreading from an infected skin wound or cellulitis; tender red streaks, fever, and enlarged regional lymph nodes may occur.

\medterm{Lymphatic Basin} The group of lymph vessels and lymph nodes that drains a particular body area. Doctors use this drainage pattern to examine nearby nodes and assess possible spread of infection or cancer.

\medterm{Lymphatic Disease} Any disorder affecting lymph vessels, lymph nodes, or related organs. It may cause swelling, repeated infection, abnormal fluid collections, or cancer and may result from birth defects, injury, infection, inflammation, or tumor growth.
\medterm{Lymphatic Filariasis} A mosquito-borne infection caused by thread-like worms that live in lymph vessels. Repeated infection can block lymph flow and cause long-term swelling of the legs, arms, breasts, or genitals.

\medterm{Lymphatic Mapping} Imaging or tracer-guided identification of lymphatic vessels and lymph nodes that drain a body area. It helps plan surgery, locate a sentinel node, investigate lymphedema, or assess possible cancer spread.

\medterm{Lymphatic Obstruction} Blocked or damaged lymph vessels or nodes that prevent normal fluid drainage. The affected area may develop persistent swelling called lymphedema.

\medterm{Lymphatic Obstruction (Lymphedema)} Alternate index label for Lymphedema.

\medterm{Lymphatic System} A network of vessels, lymph nodes, and organs that returns excess tissue fluid to the bloodstream, transports some digested fats, and helps the immune system detect and fight infection.

\medterm{Lymphatic Tissue} Lymphatic tissue contains immune cells and structures that filter lymph, trap antigens, and coordinate immune defense.

\medterm{Lymphatic Vessel} A lymphatic vessel transports lymph from tissues through lymph nodes and eventually returns fluid to the bloodstream.

\medterm{Lymphedema} Long-lasting swelling caused by poor drainage of lymph fluid. It most often affects an arm or leg and can make the skin feel tight, increase infection risk, and limit movement.

\synonyms
Lymphatic Obstruction (Lymphedema)

\medterm{Lymphedema (Lymphoedema)} Chronic swelling of a body part due to impaired lymphatic drainage, leading to
accumulation of protein-rich interstitial fluid. Primary lymphedema: caused by
congenital or hereditary lymphatic maldevelopment.

\synonyms
Lymphoedema


\medterm{Lymphedema Management} Measures used to reduce swelling and preserve function, including compression, exercise, skin care,
and selected forms of manual lymphatic therapy.

\medterm{Lymphoblast} An immature cell that normally develops into a type of infection-fighting white blood cell. Large numbers of abnormal lymphoblasts can build up in the blood or bone marrow in acute lymphoblastic leukemia.

\medterm{Lymphoblast Count} A lymphoblast count measures immature lymphoid cells in blood or bone marrow and may help evaluate or monitor acute lymphoblastic leukemia.

\medterm{Lymphocele} A lymphocele is a collection of lymphatic fluid caused by disruption or blockage of lymph vessels, often after surgery or transplantation.

\medterm{Lymphocyte} A type of white blood cell that recognizes infections and abnormal cells and helps coordinate immune defense. The main types are B cells, which make antibodies, and T cells, which direct or carry out cellular immune responses.

\medterm{Lymphocyte Activation} The process by which a lymphocyte recognizes a target and begins responding. Activated lymphocytes multiply and help coordinate antibody or cellular immune defense.
\medterm{Lymphocyte Count} A lymphocyte count measures lymphocytes in the blood and can help assess infection, immune deficiency, inflammation, or blood disorders.

\medterm{Lymphocyte Proliferation} Rapid multiplication of lymphocytes after immune stimulation. It helps build a response to infection, but abnormal or uncontrolled proliferation can occur in lymphoid cancers.
\medterm{Lymphocytic} Relating to lymphocytes or containing an increased proportion of these white blood cells. A lymphocytic pattern may occur with viral infection, chronic inflammation, autoimmune disease, or certain cancers.

\medterm{Lymphocytic Colitis} A form of microscopic colitis causing long-lasting, watery diarrhea. The colon may look normal during examination, but a biopsy shows increased immune cells in its lining; medicines, autoimmune disease, or other factors may contribute.

\synonyms
Colitis Collagenous (Lymphocytic Colitis)
Colitis Lymphocytic (Lymphocytic Colitis)
Colitis Microscopic (Lymphocytic Colitis)

\medterm{Lymphocytic Interstitial Pneumonia} A rare inflammatory lung disease in which immune cells collect in the tissue around the air sacs. It may cause cough, breathlessness, and abnormal scans and is associated with autoimmune disease or HIV.

\medterm{Lymphocytic Leukemia} Lymphocytic leukemia is a blood cancer in which abnormal lymphoid cells multiply in the bone marrow, blood, lymph nodes, or other tissues.

\medterm{Lymphocytic Mastitis} A noncancerous breast lump caused by inflammation and toughening of breast tissue. It is associated with diabetes or autoimmune disease and can look like breast cancer on examination or scans, so testing may be needed.

\medterm{Lymphocytoma Cutis} Lymphocytoma cutis is a benign skin lymphoid proliferation presenting as red or violaceous plaques or nodules and sometimes associated with infection or immune stimulation.

\medterm{Lymphocytopenia} An abnormally low number of lymphocytes, a type of white blood cell that helps fight infection. It may be temporary after illness or stress, or result from medicines, immune disorders, infections, poor marrow function, or inherited conditions. Severity and duration determine the risk of infection and need for evaluation.

\medterm{Lymphocytosis} An abnormally increased number or proportion of lymphocytes in the blood, often associated with viral infection, some bacterial infections, or lymphoproliferative disorders.

\medterm{Lymphocytosis (High Lymphocyte Count)} An increased number of lymphocytes in the blood, often associated with infection, inflammation, or a lymphoproliferative disorder.

\medterm{Lymphoedema} Chronic swelling caused by impaired lymphatic drainage, most often affecting a limb but potentially involving other regions.
\medterm{Lymphofollicular Hyperplasia} Lymphofollicular hyperplasia is enlargement or increased formation of lymphoid follicles in response to immune stimulation, infection, inflammation, or other triggers.

\medterm{Lymphogranuloma Venereum} A sexually transmitted infection caused by certain types of Chlamydia. It may begin with a small painless sore, followed by painful groin swelling or inflammation of the rectum, sometimes with bleeding or discharge.

\medterm{Lymphoid} Resembling or relating to lymphatic tissue, lymphocytes, or organs involved in lymphocyte development and immune function.

\medterm{Lymphoid Hyperplasia} Lymphoid hyperplasia is benign increase in lymphoid tissue caused by immune stimulation, infection, inflammation, or sometimes a reactive response near a tumor.

\medterm{Lymphoid Tissue} Tissue containing lymphocytes and supporting immune responses, found in lymph nodes, spleen, thymus, tonsils, and mucosal sites.
\medterm{Lymphokine} An older term for a chemical messenger released by lymphocytes, a type of white blood cell. These messengers help immune cells communicate, multiply, and coordinate responses to infection or abnormal cells.

\medterm{Lymphoma} A cancer of lymphocytes, a type of white blood cell. Hodgkin and non-Hodgkin lymphoma include many subtypes that may grow slowly or quickly and can cause swollen lymph nodes, fever, night sweats, weight loss, or organ symptoms.

\medterm{Lymphoma (Hodgkin'S And Non-Hodgkin'S)} A group of cancers arising from lymphocytes, including Hodgkin lymphoma and the
many types of non-Hodgkin lymphoma.

\synonyms
Hodgkin's and Non-Hodgkin's

\medterm{Lymphoma, Hodgkin} Alternate index label; see \textit{Hodgkin lymphoma (Hodgkin disease)}.

\medterm{Lymphoma, Non-Hodgkin} Alternate index label; see \textit{Non-Hodgkin lymphoma}.

\medterm{Lymphomatoid Papulosis} Lymphomatoid papulosis is a chronic recurrent skin disorder with cancer-like lymphoid cells; the bumps often heal spontaneously but require monitoring for lymphoma.

\medterm{Lymphopenia} Lymphopenia is an abnormally low number of lymphocytes, caused by infection, medicines, immune disease, malnutrition, marrow disorders, or other conditions.

\medterm{Lymphoplasmacytic Lymphoma} Alternate index label; see \textit{Waldenstrom macroglobulinemia}.

\medterm{Lymphopoiesis} The production and development of lymphocytes, a group of white blood cells that help fight infection and abnormal cells. These cells develop mainly in the bone marrow and thymus.

\medterm{Lymphoproliferative Disorder} A condition in which lymphocytes multiply excessively or abnormally. It may be temporary after infection or immune stimulation, or may represent a cancer such as leukemia or lymphoma.
\medterm{Lymphoscintigraphy} A scan that follows a small radioactive tracer through lymph vessels. It can map drainage to the first lymph node draining a tumor or show impaired lymph flow in lymphedema and other disorders.

\medterm{Lynch Syndrome} An inherited condition that greatly increases the risk of colorectal, womb, ovarian, stomach, urinary-tract, and several other cancers. Regular screening can detect some cancers early, and relatives may need genetic counseling or testing.

\medterm{Lysergic Acid Diethylamide} A potent hallucinogenic drug (psychedelic) classified as a Schedule 1 controlled
substance (high abuse potential, no accepted medical use). LSD acts as a partial agonist
at serotonin 5-HT2A receptors in the prefrontal cortex, leading to altered sensory
perception, mood, and cognition.

\medterm{Lysine} Lysine is an essential amino acid needed for protein production, growth, and tissue repair; it must be obtained from food.


\medterm{Lysis} The breakdown, dissolution, or destruction of cells or tissue; examples include hemolysis, the breakdown of red blood cells, and thrombolysis, the dissolution of a blood clot.

\medterm{Lysogenic Conversion} A change in a bacterium's characteristics after a virus that infects bacteria inserts its genetic material into the bacterial chromosome and provides new genes, such as genes that increase disease-causing ability.

\medterm{Lysogeny} Lysogeny is a state in which a bacteriophage's genetic material is maintained inside a bacterial cell and replicated with the host genome.

\medterm{Lysophagy} A form of selective autophagy that removes damaged lysosomes, the cell compartments that digest waste. It helps maintain cell health; failure of lysophagy can allow damaged lysosomes to accumulate and promote inflammation or cell injury.

\medterm{Lysophosphatidylcholine} A fat-like molecule found in cell membranes and blood. It helps transport fats and can influence inflammation and cell communication; abnormal levels may occur in metabolic or inflammatory diseases.

\medterm{Lysophospholipids} Lipids related to phospholipids but containing one fewer fatty-acid chain. They are parts of cell membranes and can act as signaling molecules, influencing inflammation, blood vessels, immune responses, and cell movement.

\medterm{Lysosome} A small compartment inside a cell containing substances that break down waste, worn-out cell parts, and material brought into the cell. Lysosome disorders can cause harmful material to accumulate in tissues.

\medterm{Lysozyme} A natural enzyme that helps defend against bacteria by breaking down part of their protective outer wall. It is found in tears, saliva, breast milk, and some immune cells.
\medterm{Lysozyme Test} A lysozyme test measures the enzyme lysozyme in a sample and may support evaluation of selected blood, inflammatory, or infectious disorders.


\medterm{Lyssavirus} A group of viruses that includes rabies virus and related viruses. They mainly spread through the bite or saliva of an infected mammal and can cause a severe brain infection. Prompt wound cleaning and exposure-appropriate vaccination can prevent disease before symptoms begin.

\medterm{Lytic} Causing or involving lysis, or the breakdown of cells, tissue, or another biological structure; a lytic bone lesion is an area of bone destruction.

\medterm{Laparoscopic Surgery} Surgery performed through small cuts using a camera and long instruments. It often causes less pain and faster recovery than open surgery, but still has risks such as bleeding, infection, and injury to nearby organs.

\medterm{Leber's Hereditary Optic Neuropathy} An inherited mitochondrial disorder that causes painless loss of central vision, usually first in one eye and then the other. It mainly affects young men, and genetic counseling and specialist eye care may help.

\medterm{Leigh's Syndrome} A rare progressive neurologic disorder, usually beginning in infancy or childhood, caused by problems with cellular energy production. It can damage the brainstem and cause developmental regression, movement problems, seizures, or breathing difficulties.

\medterm{Leprosy} A long-term infection caused mainly by Mycobacterium leprae that affects skin and peripheral nerves. Early treatment with multiple antibiotics prevents nerve damage and disability; casual contact does not usually spread it.

\medterm{Lewy Bodies} Abnormal deposits of the protein alpha-synuclein inside nerve cells. They are seen in Parkinson disease and dementia with Lewy bodies and are associated with movement, thinking, sleep, and visual-perception problems.

\medterm{Lichen Striatus} A usually harmless skin condition causing small, flat, pink, brown, or pale bumps in a line along a limb. It often affects children and usually clears without treatment.

\medterm{Linear Nevus Sebaceous Syndrome} A rare disorder involving a linear sebaceous nevus, a skin lesion caused by abnormal skin glands, together with possible eye, brain, or skeletal abnormalities. It requires evaluation based on the organs involved.

\medterm{Lipitor (Atorvastatin)} Lipitor is a brand name for atorvastatin, a statin medicine that lowers low-density lipoprotein cholesterol and reduces the risk of heart attack and stroke. Muscle symptoms and liver-related effects should be discussed with a clinician.

\medterm{Lowe's Syndrome} A rare inherited disorder, also called oculocerebrorenal syndrome, that mainly affects the eyes, brain, and kidneys. Features may include congenital cataracts, developmental difficulties, low muscle tone, and kidney-related loss of minerals.

\medterm{L-theanine} An amino acid found mainly in tea leaves. It may influence relaxation and alertness, but supplements can interact with medicines and should not be assumed to treat a medical condition.

\medterm{Lung Capacity} The amount of air the lungs can hold or move, measured with breathing tests such as spirometry and lung-volume testing. Results help assess conditions including asthma, chronic obstructive pulmonary disease, and restrictive lung disease.

\medterm{Lysosomal Storage Disorders} A group of inherited diseases in which missing or faulty lysosomal enzymes cause substances to build up inside cells. Accumulation can damage organs, bones, nerves, or other tissues; treatment depends on the specific disorder.

\medterm{Lytic Lesion} A lytic lesion is an area of bone destruction caused by processes such as cancer, infection, or some metabolic bone disorders.
